% Ban thu: chi giu cac cau co trong CD2

%@ id=13CD.06.007 | trang=220 | so=7 | vitri=C3.9.3,CD2.1 | muc=8 | tt=0 | gl=1
\begin{bt}
Thu gọn các biểu thức sau: $A = \dfrac{5 + \sqrt{5}}{\sqrt{5} + 2} + \dfrac{\sqrt{5}}{\sqrt{5} - 1} - \dfrac{3\sqrt{5}}{3 + \sqrt{5}}$
\loigiai{
$A = \dfrac{5 + \sqrt{5}}{\sqrt{5} + 2} + \dfrac{\sqrt{5}}{\sqrt{5} - 1} - \dfrac{3\sqrt{5}}{3 + \sqrt{5}}$\\
$= \dfrac{(5 + \sqrt{5})(\sqrt{5} - 2)}{(\sqrt{5} + 2)(\sqrt{5} - 2)} + \dfrac{\sqrt{5}(\sqrt{5} + 1)}{(\sqrt{5} - 1)(\sqrt{5} + 1)} - \dfrac{3\sqrt{5}(3 - \sqrt{5})}{(3 + \sqrt{5})(3 - \sqrt{5})}$\\
$= 3\sqrt{5} - 5 + \dfrac{5 + \sqrt{5}}{4} - \dfrac{9\sqrt{5} - 15}{4} = 3\sqrt{5} - 5 + \dfrac{5 + \sqrt{5} - 9\sqrt{5} + 15}{4}$\\
$= 3\sqrt{5} - 5 + 5 - 2\sqrt{5} = \sqrt{5}$
}
\end{bt}

%@ id=13CD.06.013 | trang=222 | so=1 | vitri=C3.9.4,CD2.1 | muc=8 | tt=0 | gl=1 | soat=ok
\begin{bt}
Rút gọn $B = \left(\dfrac{x}{x + 3\sqrt{x}} + \dfrac{1}{\sqrt{x} + 3}\right) : \left(1 - \dfrac{2}{\sqrt{x}} + \dfrac{6}{x + 3\sqrt{x}}\right)$ ($x > 0$).
\loigiai{
$B = \left(\dfrac{x}{x + 3\sqrt{x}} + \dfrac{1}{\sqrt{x} + 3}\right) : \left(1 - \dfrac{2}{\sqrt{x}} + \dfrac{6}{x + 3\sqrt{x}}\right)$ ($x > 0$)\\
$= \left(\dfrac{\sqrt{x}}{\sqrt{x} + 3} + \dfrac{1}{\sqrt{x} + 3}\right) : \left(\dfrac{\sqrt{x} - 2}{\sqrt{x}} + \dfrac{6}{\sqrt{x}(\sqrt{x} + 3)}\right)$\\
$= \dfrac{\sqrt{x} + 1}{\sqrt{x} + 3} : \dfrac{(\sqrt{x} - 2)(\sqrt{x} + 3) + 6}{\sqrt{x}(\sqrt{x} + 3)} = (\sqrt{x} + 1) \cdot \dfrac{\sqrt{x}}{x + \sqrt{x}} = 1$
}
\end{bt}

%@ id=13CD.06.014 | trang=223 | so=2 | vitri=C3.9.4,CD2.1 | muc=8 | tt=0 | gl=1
\begin{bt}
Rút gọn biểu thức $P = \dfrac{x\sqrt{2}}{2\sqrt{x} + x\sqrt{2}} + \dfrac{\sqrt{2x} - 2}{x - 2}$, với $x > 0, x \ne 2$.
\loigiai{
Với điều kiện đã cho thì:\\
$P = \dfrac{x\sqrt{2}}{\sqrt{2x}(\sqrt{2} + \sqrt{x})} + \dfrac{\sqrt{2}(\sqrt{x} - \sqrt{2})}{(\sqrt{x} - \sqrt{2})(\sqrt{x} + \sqrt{2})} = \dfrac{\sqrt{x}}{\sqrt{2} + \sqrt{x}} + \dfrac{\sqrt{2}}{\sqrt{x} + \sqrt{2}} = 1$.
}
\end{bt}

%@ id=13CD.06.023 | trang=226 | so=11 | vitri=C3.9.3,CD2.1 | muc=8 | tt=0 | gl=1
\begin{bt}
Cho biểu thức $P = \dfrac{x^3+y^3}{x^2-xy+y^2} \cdot \dfrac{x+y}{x^2-y^2}$, $x \neq \pm y$.
\begin{cacy}
\item Rút gọn biểu thức $P$.
\item Tính giá trị của $P$ khi $x = \sqrt{7-4\sqrt{3}}$ và $y = \sqrt{4-2\sqrt{3}}$.
\end{cacy}
\loigiai{
\begin{cacy}
\item $P = \dfrac{x^3+y^3}{x^2-xy+y^2} \cdot \dfrac{x+y}{x^2-y^2} = \dfrac{(x+y)(x^2-xy+y^2)}{x^2-xy+y^2} \cdot \dfrac{x+y}{(x-y)(x+y)} = \dfrac{x+y}{x-y}$.
\item Với $x = \sqrt{7-4\sqrt{3}} = 2-\sqrt{3}$ và $y = \sqrt{4-2\sqrt{3}} = \sqrt{3}-1$.

Thay vào $P$ ta được: $P = \dfrac{2-\sqrt{3} + \sqrt{3}-1}{(2-\sqrt{3}) - (\sqrt{3}-1)} = \dfrac{1}{3-2\sqrt{3}} = \dfrac{3+2\sqrt{3}}{-3}$.
\end{cacy}
}
\end{bt}

%@ id=13CD.06.025 | trang=227 | so=13 | vitri=C3.9.5,CD2.1 | muc=8 | tt=0 | gl=1
\begin{bt}
Cho biểu thức $C = \dfrac{a}{a-16} - \dfrac{2}{\sqrt{a}-4} - \dfrac{2}{\sqrt{a}+4}$.
\begin{cacy}
\item Tìm điều kiện của $a$ để biểu thức $C$ có nghĩa và rút gọn $C$.
\item Tính giá trị của biểu thức $C$ khi $a = 9-4\sqrt{5}$.
\end{cacy}
\loigiai{
\begin{cacy}
\item Biểu thức $C$ có nghĩa khi:
$\heva{&a \ge 0 \\ &a-16 \neq 0 \\ &\sqrt{a}-4 \neq 0 \\ &\sqrt{a}+4 \neq 0} \Leftrightarrow \heva{&a \ge 0 \\ &a \neq 16 \\ &a \neq 16 \\ &\forall a \ge 0} \Rightarrow a \ge 0, a \neq 16$.

Rút gọn $C = \dfrac{a}{a-16} - \dfrac{2}{\sqrt{a}-4} - \dfrac{2}{\sqrt{a}+4} = \dfrac{a}{\left(\sqrt{a}-4\right)\left(\sqrt{a}+4\right)} - \dfrac{2}{\sqrt{a}-4} - \dfrac{2}{\sqrt{a}+4}$
\[
\begin{aligned}
&= \dfrac{a - 2\left(\sqrt{a}+4\right) - 2\left(\sqrt{a}-4\right)}{\left(\sqrt{a}+4\right)\left(\sqrt{a}-4\right)} = \dfrac{a - 2\sqrt{a} - 8 - 2\sqrt{a} + 8}{\left(\sqrt{a}+4\right)\left(\sqrt{a}-4\right)} \\
&= \dfrac{a-4\sqrt{a}}{\left(\sqrt{a}+4\right)\left(\sqrt{a}-4\right)} = \dfrac{\sqrt{a}\left(\sqrt{a}-4\right)}{\left(\sqrt{a}-4\right)\left(\sqrt{a}+4\right)} = \dfrac{\sqrt{a}}{\sqrt{a}+4}.
\end{aligned}
\]

\item Giá trị của $C$ khi $a = 9-4\sqrt{5}$.
Ta có: $a = 9-4\sqrt{5} = 4 - 4\sqrt{5} + 5 = \left(2-\sqrt{5}\right)^2 \Rightarrow \sqrt{a} = \sqrt{\left(2-\sqrt{5}\right)^2} = \sqrt{5}-2$.

Vậy $C = \dfrac{\sqrt{a}}{\sqrt{a}+4} = \dfrac{\sqrt{5}-2}{\sqrt{5}-2+4} = \dfrac{\sqrt{5}-2}{\sqrt{5}+2} = 9-4\sqrt{5}$.
\end{cacy}
}
\end{bt}

%@ id=13CD.06.028 | trang=229 | so=16 | vitri=C3.9.4,CD2.1 | muc=8 | tt=0 | ghichu=đề và đáp số đã đối chiếu trang gốc 229 (dấu nhân) | gl=1 | soat=sua
\begin{bt}
Cho biểu thức $P = \dfrac{\left(\sqrt{a}-\sqrt{b}\right)^2 + 4\sqrt{ab}}{\sqrt{a}+\sqrt{b}} \cdot \dfrac{a\sqrt{b}-b\sqrt{a}}{\sqrt{ab}}$.
\begin{cacy}
\item Tìm điều kiện để $P$ có nghĩa
\item Rút gọn $P$
\item Tính giá trị của $P$ khi $a = 2\sqrt{3}; b = \sqrt{3}$
\end{cacy}
\loigiai{
\begin{cacy}
\item Để $P$ có nghĩa thì $\heva{&a \ge 0 \\ &b \ge 0 \\ &ab \ge 0 \\ &\sqrt{a}+\sqrt{b} \neq 0 \\ &\sqrt{ab} \neq 0} \Leftrightarrow \heva{&a > 0 \\ &b > 0}$.
Vậy $P$ có nghĩa khi $a > 0$ và $b > 0$.

\item Với $a > 0$ và $b > 0$ ta có: $P = \dfrac{\left(\sqrt{a}-\sqrt{b}\right)^2 + 4\sqrt{ab}}{\sqrt{a}+\sqrt{b}} \cdot \dfrac{a\sqrt{b}-b\sqrt{a}}{\sqrt{ab}} = a - b$.
Vậy $P = a - b$ ($a > 0; b > 0$).

\item Ta có $\heva{&a = 2\sqrt{3} \\ &b = \sqrt{3}}$ (thỏa mãn điều kiện) thay vào biểu thức $P$ ta được $P = 2\sqrt{3} - \sqrt{3} = \sqrt{3}$.
Vậy $P = \sqrt{3}$ khi $a = 2\sqrt{3}; b = \sqrt{3}$.
\end{cacy}
}
\end{bt}

%@ id=13CD.06.029 | trang=229 | so=17 | vitri=C3.9.4,CD2.1 | muc=8 | tt=0 | gl=1
\begin{bt}
Cho biểu thức $A = \left( \dfrac{\sqrt{a}}{2} - \dfrac{1}{2\sqrt{a}} \right) \left( \dfrac{a-\sqrt{a}}{\sqrt{a}+1} - \dfrac{a+\sqrt{a}}{\sqrt{a}-1} \right)$.
\begin{cacy}
\item Rút gọn biểu thức $A$
\item Tính giá trị của $A$ khi $|a-1| = 1$
\end{cacy}
\loigiai{
\begin{cacy}
\item Điều kiện để $A$ có nghĩa là: $a > 0; a \neq 1$.
$A = \left( \dfrac{\sqrt{a}}{2} - \dfrac{1}{2\sqrt{a}} \right) \left( \dfrac{a-\sqrt{a}}{\sqrt{a}+1} - \dfrac{a+\sqrt{a}}{\sqrt{a}-1} \right) = \left( \dfrac{a-1}{2\sqrt{a}} \right) \left( \dfrac{-4a}{a-1} \right) = -2\sqrt{a}$.
Vậy $A = -2\sqrt{a}$ ($a > 0; a \neq 1$).

\item Ta có $|a-1| = 1 \Leftrightarrow \hoac{&a-1 = 1 \\ &a-1 = -1} \Leftrightarrow \hoac{&a = 2 \text{ (tm)} \\ &a = 0 \text{ (loại)}}$.
Thay $a = 2$ vào biểu thức $A$ ta được $A = -2\sqrt{2}$.
\end{cacy}
}
\end{bt}

%@ id=13CD.06.031 | trang=230 | so=19 | vitri=C3.9.4,CD2.1 | muc=8 | tt=0 | gl=1
\begin{bt}
Cho ba biểu thức $M = \dfrac{x\sqrt{x}-8}{3+\left(\sqrt{x}+1\right)^2}$; $N = \dfrac{\left(\sqrt{x}+1\right)^3 - \left(\sqrt{x}-1\right)^3}{(x-4)(3x+1)}$; $P = \dfrac{\sqrt{x}}{2+\sqrt{x}}$.
\begin{cacy}
\item Tìm tất cả các số thực $x$ thỏa mãn $M = x-4$.
\item Trong trường hợp các biểu thức $M, N$ và $P$ xác định, rút gọn biểu thức $Q = MN + P$.
\end{cacy}
\loigiai{
\begin{cacy}
\item Ta có $M = \dfrac{x\sqrt{x}-8}{3+\left(\sqrt{x}+1\right)^2} = \dfrac{\left(\sqrt{x}\right)^3 - 2^3}{3+\left(\sqrt{x}+1\right)^2} = \dfrac{\left(\sqrt{x}-2\right)\left(x+2\sqrt{x}+4\right)}{3 + x + 2\sqrt{x} + 1} = \dfrac{\left(\sqrt{x}-2\right)\left(x+2\sqrt{x}+4\right)}{x+2\sqrt{x}+4} = \sqrt{x}-2$.

Khi đó $M = x-4 \Leftrightarrow \sqrt{x}-2 = x-4 = \left(\sqrt{x}\right)^2 - 2^2 \Leftrightarrow \sqrt{x}-2 = \left(\sqrt{x}-2\right)\left(\sqrt{x}+2\right)$
$\Leftrightarrow \left(\sqrt{x}-2\right)\left(\sqrt{x}+2-1\right) = 0 \Leftrightarrow \hoac{&\sqrt{x}-2 = 0 \\ &\sqrt{x}+1 = 0} \Leftrightarrow x = 4$ (thỏa mãn).
Vậy $x = 4$ thì $M = x-4$.

\item Điều kiện $x \ge 0; x \neq 4$.
Ta có $M = \sqrt{x}-2$; $P = \dfrac{\sqrt{x}}{2+\sqrt{x}}$.
\[
\begin{aligned}
N &= \dfrac{\left(\sqrt{x}+1\right)^3 - \left(\sqrt{x}-1\right)^3}{(x-4)(3x+1)} = \dfrac{\left(\sqrt{x}+1-\sqrt{x}+1\right) \left[\left(\sqrt{x}+1\right)^2 + \left(\sqrt{x}+1\right)\left(\sqrt{x}-1\right) + \left(\sqrt{x}-1\right)^2\right]}{(x-4)(3x+1)} \\
&= \dfrac{2\left(x+2\sqrt{x}+1 + x-1 + x-2\sqrt{x}+1\right)}{(x-4)(3x+1)} = \dfrac{2(3x+1)}{(x-4)(3x+1)} = \dfrac{2}{x-4}.
\end{aligned}
\]
$\Rightarrow Q = M \cdot N + P = \left(\sqrt{x}-2\right) \cdot \dfrac{2}{x-4} + \dfrac{\sqrt{x}}{\sqrt{x}+2} = \dfrac{2}{\sqrt{x}+2} + \dfrac{\sqrt{x}}{\sqrt{x}+2} = 1$.
Vậy $Q = 1$.
\end{cacy}
}
\end{bt}

%@ id=13CD.06.049 | trang=241 | so=37 | vitri=C3.9.4,CD2.1 | muc=8 | tt=0 | gl=1
\begin{bt}
Cho biểu thức $P = \left( 1 - \dfrac{1}{\sqrt{x}} \right) : \left( \dfrac{\sqrt{x}-1}{\sqrt{x}} + \dfrac{1-\sqrt{x}}{x+\sqrt{x}} \right)$, (với $x > 0$ và $x \ne 1$).
\begin{cacy}
\item Rút gọn biểu thức $P$.
\item Tính giá trị của biểu thức $P$ tại $x = \sqrt{2022+4\sqrt{2018}} - \sqrt{2022-4\sqrt{2018}}$.
\end{cacy}
\loigiai{
\begin{cacy}
\item Ta có $1 - \dfrac{1}{\sqrt{x}} = \dfrac{\sqrt{x}-1}{\sqrt{x}}$\\
Và $\dfrac{\sqrt{x}-1}{\sqrt{x}} + \dfrac{1-\sqrt{x}}{x+\sqrt{x}} = \dfrac{x-1+1-\sqrt{x}}{\sqrt{x}\left(\sqrt{x}+1\right)} = \dfrac{\left(\sqrt{x}-1\right)\sqrt{x}}{\left(\sqrt{x}+1\right)\sqrt{x}} = \dfrac{\sqrt{x}-1}{\sqrt{x}+1}$\\
nên $P = \dfrac{\sqrt{x}-1}{\sqrt{x}} \cdot \dfrac{\sqrt{x}+1}{\sqrt{x}-1} = \dfrac{\sqrt{x}+1}{\sqrt{x}}$.

\item Có $x = \sqrt{2022+4\sqrt{2018}} - \sqrt{2022-4\sqrt{2018}} = \sqrt{\left(\sqrt{2018}+2\right)^2} - \sqrt{\left(\sqrt{2018}-2\right)^2}$\\
$= \left|\sqrt{2018}+2\right| - \left|\sqrt{2018}-2\right| = \sqrt{2018}+2 - \sqrt{2018}+2 = 4$ thỏa mãn điều kiện $x > 0$ và $x \ne 1$.\\
Vậy giá trị của biểu thức $P$ tại $x = 4$ là: $\dfrac{\sqrt{4}+1}{\sqrt{4}} = \dfrac{3}{2}$.
\end{cacy}
}
\end{bt}

%@ id=13CD.06.008 | trang=220 | so=8 | vitri=C3.8.3,CD2.1 | muc=9 | tt=0 | gl=1
\begin{bt}
Tính $B = 21\left(\sqrt{2 + \sqrt{3}} + \sqrt{3 - \sqrt{5}}\right)^2 - 6\left(\sqrt{2 - \sqrt{3}} + \sqrt{3 + \sqrt{5}}\right)^2 - 15\sqrt{15}$.
\loigiai{
$B = \dfrac{21}{2}\left(\sqrt{4 + 2\sqrt{3}} + \sqrt{6 - 2\sqrt{5}}\right)^2 - 3\left(\sqrt{4 - 2\sqrt{3}} + \sqrt{6 + 2\sqrt{5}}\right)^2 - 15\sqrt{15}$\\
$= \dfrac{21}{2}\left(\sqrt{3} + 1 + \sqrt{5} - 1\right)^2 - 3\left(\sqrt{3} - 1 + \sqrt{5} + 1\right)^2 - 15\sqrt{15}$\\
$= \dfrac{15}{2}\left(\sqrt{3} + \sqrt{5}\right)^2 - 15\sqrt{15} = 60$.
}
\end{bt}

\dang{Tìm điều kiện xác định của biểu thức}

\begin{kienthuc}
\textbf{1. Phương pháp}

\begin{tabular}{|c|l|c|l|}\hline
\multicolumn{2}{|c|}{\textbf{BIỂU THỨC - ĐKXĐ:}} & \multicolumn{2}{c|}{\textbf{CÂU}} \\ \hline
\textbf{1.} & $\sqrt{A}$ \quad ĐKXĐ: $A \ge 0$ & Câu: $\sqrt{x - 2018}$ & ĐKXĐ: $x \ge 2018$ \\ \hline
\textbf{2.} & $\dfrac{A}{B}$ \quad ĐKXĐ: $B \ne 0$ & Câu: $\dfrac{x + 4}{x - 7}$ & ĐKXĐ: $x \ne 7$ \\ \hline
\textbf{3.} & $\dfrac{A}{\sqrt{B}}$ \quad ĐKXĐ: $B > 0$ & Câu: $\dfrac{x + 1}{\sqrt{x - 3}}$ & ĐKXĐ: $x > 3$ \\ \hline
\textbf{4.} & $\dfrac{\sqrt{A}}{\sqrt{B}}$ \quad ĐKXĐ: $A \ge 0; B > 0$ & Câu: $\dfrac{\sqrt{x}}{\sqrt{x - 3}}$ & ĐKXĐ: $\heva{&x \ge 0\\ &x > 3} \Leftrightarrow x > 3$ \\ \hline
\textbf{5.} & $\sqrt{\dfrac{A}{B}}$ \quad ĐKXĐ: $\hoac{\&\heva{&A \le 0\\ &B < 0}\\ &\heva{&A \ge 0\\ &B > 0}}$ & Câu: $\sqrt{\dfrac{x + 1}{x + 2}}$ & ĐKXĐ: $\hoac{\&\heva{&x + 1 \le 0\\ &x + 2 < 0}\\ &\heva{&x + 1 \ge 0\\ &x + 2 > 0}} \Leftrightarrow \hoac{&x < -2\\ &x \ge -1}$ \\ \hline
\textbf{6.} & \begin{tabular}[t]{l} Cho $a > 0$ ta có:\\ $x^2 > a \Leftrightarrow \hoac{&x > \sqrt{a}\\ &x < -\sqrt{a}}$ \end{tabular} & Câu: $x^2 > 1$ & $\Leftrightarrow \hoac{&x > \sqrt{a}\\ &x < -\sqrt{a}}$ \\ \hline
\textbf{7.} & \begin{tabular}[t]{l} Cho $a > 0$ ta có:\\ $x^2 < a \Leftrightarrow -\sqrt{a} < x < \sqrt{a}$ \end{tabular} & Câu: $x^2 < 4$ & $\Leftrightarrow -2 < x < 2$ \\ \hline
\end{tabular}
\end{kienthuc}

%@ id=13CD.06.016 | trang=223 | so=4 | vitri=C3.9.4,CD2.1 | muc=9 | tt=0 | ghichu=đã đối chiếu trang gốc, khớp gốc, toán đúng | gl=1 | soat=sua
\begin{bt}
Rút gọn các biểu thức:
\begin{cacy}
\item $A = \sqrt{x} - \sqrt{x - \sqrt{x} + \dfrac{1}{4}}$ khi $x \ge 0$.
\item $B = \sqrt{4x - 2\sqrt{4x - 1}} + \sqrt{4x + 2\sqrt{4x - 1}}$ khi $x \ge \dfrac{1}{4}$.
\end{cacy}
\loigiai{
a) $A = \sqrt{x} - \sqrt{x - \sqrt{x} + \dfrac{1}{4}} = \sqrt{x} - \sqrt{\left(\sqrt{x} - \dfrac{1}{2}\right)^2} = \sqrt{x} - \left|\sqrt{x} - \dfrac{1}{2}\right|$\\
+ Nếu $\sqrt{x} \ge \dfrac{1}{2} \Leftrightarrow x \ge \dfrac{1}{4}$ thì $\left|\sqrt{x} - \dfrac{1}{2}\right| = \sqrt{x} - \dfrac{1}{2} \Rightarrow A = \dfrac{1}{2}$.\\
+ Nếu $\sqrt{x} < \dfrac{1}{2} \Leftrightarrow 0 \le x < \dfrac{1}{4}$ thì $\left|\sqrt{x} - \dfrac{1}{2}\right| = -\sqrt{x} + \dfrac{1}{2} \Rightarrow A = 2\sqrt{x} - \dfrac{1}{2}$\\
b) $B = \sqrt{4x - 2\sqrt{4x - 1}} + \sqrt{4x + 2\sqrt{4x - 1}} = \sqrt{4x - 1 - 2\sqrt{4x - 1} + 1} + \sqrt{4x - 1 + 2\sqrt{4x - 1} + 1}$\\
Hay $B = \sqrt{\left(\sqrt{4x - 1} - 1\right)^2} + \sqrt{\left(\sqrt{4x - 1} + 1\right)^2} = \left|\sqrt{4x - 1} - 1\right| + \left|\sqrt{4x - 1} + 1\right| = \left|\sqrt{4x - 1} - 1\right| + \sqrt{4x - 1} + 1$\\
+ Nếu $\sqrt{4x - 1} - 1 \ge 0 \Leftrightarrow 4x - 1 \ge 1 \Leftrightarrow x \ge \dfrac{1}{2}$ thì $\left|\sqrt{4x - 1} - 1\right| = \sqrt{4x - 1} - 1$ suy ra $B = 2\sqrt{4x - 1}$.\\
+ Nếu $\sqrt{4x - 1} - 1 < 0 \Leftrightarrow 4x - 1 < 1 \Leftrightarrow \dfrac{1}{4} \le x < \dfrac{1}{2}$ thì $\left|\sqrt{4x - 1} - 1\right| = -\sqrt{4x - 1} + 1$ suy ra $B = 2$.
}
\end{bt}

%@ id=13CD.06.017 | trang=224 | so=5 | vitri=C3.9.4,CD2.1 | muc=9 | tt=0 | ghichu=gốc nhầm mẫu thành (\sqrt{a}+\sqrt{b})^3 ở dòng 1 lời giải và thừa = 0 ở dòng 2 | gl=1 | soat=sua
\begin{bt}
Cho các số thực dương $a, b; a \ne b$.
Chứng minh rằng:
$$\dfrac{\dfrac{(a - b)^3}{\left(\sqrt{a} - \sqrt{b}\right)^3} - b\sqrt{b} + 2a\sqrt{a}}{a\sqrt{a} - b\sqrt{b}} + \dfrac{3a + 3\sqrt{ab}}{b - a} = 0.$$
\loigiai{
Ta có: $Q = \dfrac{\dfrac{(a - b)^3}{\left(\sqrt{a} - \sqrt{b}\right)^3} - b\sqrt{b} + 2a\sqrt{a}}{a\sqrt{a} - b\sqrt{b}} + \dfrac{3a + 3\sqrt{ab}}{b - a}$\\
$= \dfrac{\dfrac{\left(\sqrt{a} - \sqrt{b}\right)^3 \left(\sqrt{a} + \sqrt{b}\right)^3}{\left(\sqrt{a} - \sqrt{b}\right)^3} - b\sqrt{b} + 2a\sqrt{a}}{\left(\sqrt{a} - \sqrt{b}\right)(a + \sqrt{ab} + b)} - \dfrac{3\sqrt{a}\left(\sqrt{a} + \sqrt{b}\right)}{\left(\sqrt{a} - \sqrt{b}\right)\left(\sqrt{a} + \sqrt{b}\right)}$\\
$= \dfrac{a\sqrt{a} + 3a\sqrt{b} + 3b\sqrt{a} + b\sqrt{b} - b\sqrt{b} + 2a\sqrt{a}}{\left(\sqrt{a} - \sqrt{b}\right)(a + \sqrt{ab} + b)} - \dfrac{3\sqrt{a}}{\left(\sqrt{a} - \sqrt{b}\right)}$\\
$= \dfrac{3a\sqrt{a} + 3a\sqrt{b} + 3b\sqrt{a} - 3a\sqrt{a} - 3a\sqrt{b} - 3b\sqrt{a}}{\left(\sqrt{a} - \sqrt{b}\right)(a + \sqrt{ab} + b)} = 0$
}
\end{bt}

%@ id=13CD.06.018 | trang=224 | so=6 | vitri=C3.9.4,CD2.1 | muc=9 | tt=0 | ghichu=gốc biến đổi sai ở dòng 2 lời giải | gl=1 | soat=sua
\begin{bt}
Rút gọn biểu thức $A = \dfrac{x + \sqrt{x} - 6}{x - 9} + \dfrac{x - 7\sqrt{x} + 19}{x + \sqrt{x} - 12} - \dfrac{x - 5\sqrt{x}}{x + 4\sqrt{x}}$; $x > 0, x \ne 9$.
\loigiai{
$A = \dfrac{x + \sqrt{x} - 6}{x - 9} + \dfrac{x - 7\sqrt{x} + 19}{x + \sqrt{x} - 12} - \dfrac{x - 5\sqrt{x}}{x + 4\sqrt{x}} = \dfrac{\sqrt{x} - 2}{\sqrt{x} - 3} + \dfrac{x - 7\sqrt{x} + 19}{(\sqrt{x} - 3)(\sqrt{x} + 4)} - \dfrac{\sqrt{x} - 5}{\sqrt{x} + 4}$\\
$= \dfrac{x + 2\sqrt{x} - 8 + x - 7\sqrt{x} + 19 - x + 8\sqrt{x} - 15}{(\sqrt{x} - 3)(\sqrt{x} + 4)} = \dfrac{(\sqrt{x} - 1)(\sqrt{x} + 4)}{(\sqrt{x} - 3)(\sqrt{x} + 4)} = \dfrac{\sqrt{x} - 1}{\sqrt{x} - 3}$.
}
\end{bt}

%@ id=13CD.06.027 | trang=228 | so=15 | vitri=C3.9.4,CD2.2 | muc=7 | tt=0 | gl=1 | soat=ok
\begin{bt}
Cho biểu thức $P = \left( \dfrac{x-2}{x+2\sqrt{x}} + \dfrac{1}{\sqrt{x}+2} \right) \cdot \dfrac{\sqrt{x}+1}{\sqrt{x}-1}$ với $x > 0$ và $x \neq 1$.
\begin{cacy}
\item Chứng minh rằng $P = \dfrac{\sqrt{x}+1}{\sqrt{x}}$.
\item Tìm các giá trị của $x$ để $2P = 2\sqrt{x} + 5$.
\end{cacy}
\loigiai{
\begin{cacy}
\item $P = \left( \dfrac{x-2+\sqrt{x}}{\sqrt{x}\left(\sqrt{x}+2\right)} \right) \cdot \dfrac{\sqrt{x}+1}{\sqrt{x}-1} = \dfrac{\left(\sqrt{x}-1\right)\left(\sqrt{x}+2\right)}{\sqrt{x}\left(\sqrt{x}+2\right)} \cdot \dfrac{\sqrt{x}+1}{\sqrt{x}-1} = \dfrac{\sqrt{x}+1}{\sqrt{x}}$.

\item Theo câu a) $P = \dfrac{\sqrt{x}+1}{\sqrt{x}}$
$\Rightarrow 2P = 2\sqrt{x}+5 \Leftrightarrow \dfrac{2\sqrt{x}+2}{\sqrt{x}} = 2\sqrt{x}+5 \Leftrightarrow 2\sqrt{x}+2 = 2x+5\sqrt{x} \Leftrightarrow 2x+3\sqrt{x}-2 = 0$ và $x > 0$
$\Leftrightarrow \left(\sqrt{x}+2\right)\left(\sqrt{x}-\dfrac{1}{2}\right) = 0 \Leftrightarrow \sqrt{x} = \dfrac{1}{2} \Leftrightarrow x = \dfrac{1}{4}$.
\end{cacy}
}
\end{bt}

%@ id=13CD.06.019 | trang=224 | so=7 | vitri=C3.9.5,CD2.2 | muc=8 | tt=0 | gl=1 | soat=ok
\begin{bt}
Cho biểu thức $P = \dfrac{3x + 5\sqrt{x} - 4}{(\sqrt{x} + 3)(\sqrt{x} - 1)} - \dfrac{\sqrt{x} + 1}{\sqrt{x} + 3} - \dfrac{\sqrt{x} + 3}{\sqrt{x} - 1}$ (với $x \ge 0; x \ne 1$)
\begin{cacy}
\item Rút gọn biểu thức $P$.
\item Tìm $x$ sao cho $P = \dfrac{1}{2}$.
\end{cacy}
\loigiai{
a) Với $x \ge 0; x \ne 1$ ta có:\\
$P = \dfrac{3x + 5\sqrt{x} - 4}{(\sqrt{x} + 3)(\sqrt{x} - 1)} - \dfrac{\sqrt{x} + 1}{\sqrt{x} + 3} - \dfrac{\sqrt{x} + 3}{\sqrt{x} - 1}$\\
$= \dfrac{3x + 5\sqrt{x} - 4 - (\sqrt{x} + 1)(\sqrt{x} - 1) - (\sqrt{x} + 3)^2}{(\sqrt{x} + 3)(\sqrt{x} - 1)}$\\
$= \dfrac{3x + 5\sqrt{x} - 4 - (x - 1) - (x + 6\sqrt{x} + 9)}{(\sqrt{x} + 3)(\sqrt{x} - 1)} = \dfrac{x - \sqrt{x} - 12}{(\sqrt{x} + 3)(\sqrt{x} - 1)} = \dfrac{x - 4\sqrt{x} + 3\sqrt{x} - 12}{(\sqrt{x} + 3)(\sqrt{x} - 1)}$\\
$= \dfrac{\sqrt{x}(\sqrt{x} - 4) + 3(\sqrt{x} - 4)}{(\sqrt{x} + 3)(\sqrt{x} - 1)} = \dfrac{(\sqrt{x} - 4)(\sqrt{x} + 3)}{(\sqrt{x} + 3)(\sqrt{x} - 1)} = \dfrac{\sqrt{x} - 4}{\sqrt{x} - 1}$.\\
b) $P = \dfrac{1}{2} \Leftrightarrow \dfrac{\sqrt{x} - 4}{\sqrt{x} - 1} = \dfrac{1}{2} \Leftrightarrow 2\sqrt{x} - 8 = \sqrt{x} - 1 \Leftrightarrow \sqrt{x} = 7 \Leftrightarrow x = 49$ (thoả điều kiện)\\
Vậy khi $x = 49$ thì $P = \dfrac{1}{2}$.
}
\end{bt}

%% doan 225-234

%@ id=13CD.06.024 | trang=227 | so=12 | vitri=C3.9.5,CD2.2 | muc=8 | tt=0 | gl=1
\begin{bt}
Cho biểu thức $A = \dfrac{1}{2+\sqrt{x}} + \dfrac{1}{2-\sqrt{x}} - \dfrac{2\sqrt{x}}{4-x}$ ($x \ge 0, x \neq 4$). Rút gọn $A$ và tìm $x$ để $A = \dfrac{1}{3}$.
\loigiai{
$A = \dfrac{1}{2+\sqrt{x}} + \dfrac{1}{2-\sqrt{x}} - \dfrac{2\sqrt{x}}{4-x} = \dfrac{4}{4-x} - \dfrac{2\sqrt{x}}{4-x} = \dfrac{2\left(2-\sqrt{x}\right)}{4-x} = \dfrac{2}{2+\sqrt{x}}$.

Với $A = \dfrac{1}{3} \Leftrightarrow \dfrac{2}{2+\sqrt{x}} = \dfrac{1}{3} \Leftrightarrow \sqrt{x} = 4 \Leftrightarrow x = 16$ (nhận).

Vậy $A = \dfrac{1}{3}$ khi $x = 16$.
}
\end{bt}

%@ id=13CD.06.050 | trang=241 | so=38 | vitri=C3.9.4,CD2.2 | muc=8 | tt=0 | ghichu=lời giải câu b gốc in nhầm mẫu x-15 | gl=1 | soat=sua
\begin{bt}
Cho hai biểu thức $A = \dfrac{\sqrt{x}+2}{\sqrt{x}-5}$ và $B = \dfrac{3}{\sqrt{x}+5} + \dfrac{20-2\sqrt{x}}{x-25}$ với $x \ge 0, x \ne 25$.
\begin{cacy}
\item Tính giá trị biểu thức $A$ khi $x = 9$.
\item Chứng minh rằng $B = \dfrac{1}{\sqrt{x}-5}$.
\item Tìm tất cả các giá trị của $x$ để $A = B \cdot |x-4|$.
\end{cacy}
\loigiai{
\begin{cacy}
\item Khi $x = 9$ ta có $A = \dfrac{\sqrt{9}+2}{\sqrt{9}-5} = \dfrac{3+2}{3-5} = -\dfrac{5}{2}$.

\item Với $x \ge 0, x \ne 25$ thì
\begin{align*}
B &= \dfrac{3}{\sqrt{x}+5} + \dfrac{20-2\sqrt{x}}{x-25} = \dfrac{3}{\sqrt{x}+5} + \dfrac{20-2\sqrt{x}}{\left(\sqrt{x}+5\right)\left(\sqrt{x}-5\right)} \\
&= \dfrac{3\left(\sqrt{x}-5\right)+20-2\sqrt{x}}{\left(\sqrt{x}+5\right)\left(\sqrt{x}-5\right)} = \dfrac{3\sqrt{x}-15+20-2\sqrt{x}}{\left(\sqrt{x}+5\right)\left(\sqrt{x}-5\right)} \\
&= \dfrac{\sqrt{x}+5}{\left(\sqrt{x}+5\right)\left(\sqrt{x}-5\right)} = \dfrac{1}{\sqrt{x}-5} \quad \text{(đpcm)}.
\end{align*}

\item Với $x \ge 0, x \ne 25$ Ta có: $A = B \cdot |x-4|$\\
$\Leftrightarrow \dfrac{\sqrt{x}+2}{\sqrt{x}-5} = \dfrac{1}{\sqrt{x}-5} \cdot |x-4| \Leftrightarrow \sqrt{x}+2 = |x-4| \quad (*)$\\
Nếu $x \ge 4, x \ne 25$ thì $(*)$ trở thành: $\sqrt{x}+2 = x-4$\\
$\Leftrightarrow x-\sqrt{x}-6 = 0 \Leftrightarrow \left(\sqrt{x}-3\right)\left(\sqrt{x}+2\right) = 0$\\
Do $\sqrt{x}+2 > 0$ nên $\sqrt{x} = 3 \Leftrightarrow x = 9$ (thỏa mãn).\\
Nếu $0 \le x < 4$ thì $(*)$ trở thành: $\sqrt{x}+2 = 4-x$\\
$\Leftrightarrow x+\sqrt{x}-2 = 0 \Leftrightarrow \left(\sqrt{x}-1\right)\left(\sqrt{x}+2\right) = 0$\\
Do $\sqrt{x}+2 > 0$ nên $\sqrt{x} = 1 \Leftrightarrow x = 1$ (thỏa mãn).\\
Vậy có hai giá trị $x = 1$ và $x = 9$ thỏa mãn yêu cầu bài toán.
\end{cacy}
}
\end{bt}

%@ id=13CD.06.053 | trang=245 | so=41 | vitri=C3.9.4,CD2.2 | muc=8 | tt=0 | ghichu=lời giải câu b gốc đổi nhầm tử số, câu c ghi \sqrt{x}+1>0 | gl=1 | soat=sua
\begin{bt}
Cho biểu thức $P = \left( 1 - \dfrac{\sqrt{x}}{\sqrt{x}+1} \right) : \left( \dfrac{\sqrt{x}+2}{\sqrt{x}+3} + \dfrac{\sqrt{x}-3}{2-\sqrt{x}} + \dfrac{\sqrt{x}-2}{x+\sqrt{x}-6} \right)$ với $x \ge 0; x \ne 4$.
\begin{cacy}
\item Rút gọn $P$
\item Tính giá trị của $P$ biết $x = \dfrac{3-\sqrt{5}}{2}$
\item Tìm $x \in \mathbb{Z}$ để $P \in \mathbb{Z}$
\item So sánh $P$ với $1$
\item Tìm các giá trị của $x$ để $P = \sqrt{x}-3$.
\end{cacy}
\loigiai{
a) Với $x \ge 0; x \ne 4$:
\[
\begin{aligned}
P &= \left( 1 - \dfrac{\sqrt{x}}{\sqrt{x}+1} \right) : \left( \dfrac{\sqrt{x}+2}{\sqrt{x}+3} + \dfrac{\sqrt{x}-3}{2-\sqrt{x}} + \dfrac{\sqrt{x}-2}{x+\sqrt{x}-6} \right) \\
&= \dfrac{\sqrt{x}+1-\sqrt{x}}{\sqrt{x}+1} : \left( \dfrac{\sqrt{x}+2}{\sqrt{x}+3} + \dfrac{\sqrt{x}-3}{2-\sqrt{x}} + \dfrac{\sqrt{x}-2}{x+3\sqrt{x}-2\sqrt{x}-6} \right) \\
&= \dfrac{1}{\sqrt{x}+1} : \left[ \dfrac{\sqrt{x}+2}{\sqrt{x}+3} + \dfrac{\sqrt{x}-3}{2-\sqrt{x}} + \dfrac{\sqrt{x}-2}{\sqrt{x}(\sqrt{x}+3)-2(\sqrt{x}+3)} \right] \\
&= \dfrac{1}{\sqrt{x}+1} : \left[ \dfrac{\sqrt{x}+2}{\sqrt{x}+3} - \dfrac{\sqrt{x}-3}{\sqrt{x}-2} + \dfrac{\sqrt{x}-2}{(\sqrt{x}+3)(\sqrt{x}-2)} \right] \\
&= \dfrac{1}{\sqrt{x}+1} : \dfrac{(\sqrt{x}+2)(\sqrt{x}-2) - (\sqrt{x}-3)(\sqrt{x}+3) + \sqrt{x}-2}{(\sqrt{x}+3)(\sqrt{x}-2)} \\
&= \dfrac{1}{\sqrt{x}+1} : \dfrac{x-4-x+9+\sqrt{x}-2}{(\sqrt{x}+3)(\sqrt{x}-2)} = \dfrac{1}{\sqrt{x}+1} \cdot \dfrac{(\sqrt{x}+3)(\sqrt{x}-2)}{\sqrt{x}+3} = \dfrac{\sqrt{x}-2}{\sqrt{x}+1}.
\end{aligned}
\]
Vậy với $x \ge 0; x \ne 4$ thì $P = \dfrac{\sqrt{x}-2}{\sqrt{x}+1}$.

b) Với $x = \dfrac{3-\sqrt{5}}{2} = \dfrac{6-2\sqrt{5}}{4} = \dfrac{5-2\sqrt{5}\cdot 1 + 1}{4} = \dfrac{(\sqrt{5}-1)^2}{4}$ thỏa mãn $x \ge 0; x \ne 4$
\[
\Rightarrow \sqrt{x} = \sqrt{\dfrac{(\sqrt{5}-1)^2}{4}} = \dfrac{|\sqrt{5}-1|}{2} = \dfrac{\sqrt{5}-1}{2} \text{ (vì } \sqrt{5} > 1)
\]
\[
\Rightarrow P = \dfrac{\dfrac{\sqrt{5}-1}{2}-2}{\dfrac{\sqrt{5}-1}{2}+1} = \dfrac{\dfrac{\sqrt{5}-5}{2}}{\dfrac{\sqrt{5}+1}{2}} = \dfrac{\sqrt{5}-5}{\sqrt{5}+1} = \dfrac{(\sqrt{5}-5)(\sqrt{5}-1)}{(\sqrt{5}+1)(\sqrt{5}-1)}
\]
\[
= \dfrac{5-\sqrt{5}-5\sqrt{5}+5}{4} = \dfrac{10-6\sqrt{5}}{4} = \dfrac{5-3\sqrt{5}}{2}.
\]

c) $P = \dfrac{\sqrt{x}-2}{\sqrt{x}+1} = \dfrac{\sqrt{x}+1-3}{\sqrt{x}+1} = 1 - \dfrac{3}{\sqrt{x}+1}$

$P \in \mathbb{Z} \Leftrightarrow \dfrac{3}{\sqrt{x}+1} \in \mathbb{Z} \Leftrightarrow 3 \vdots (\sqrt{x}+1) \Leftrightarrow \sqrt{x}+1$ là $\text{Ư}(3)$

$\text{Ư}(3) = \{\pm 1; \pm 3\}$

Mà $\sqrt{x} \ge 0, \forall x \ge 0; x \ne 4 \Rightarrow \sqrt{x}+1 > 0, \forall x \ge 0; x \ne 4$ nên:

TH1: $\sqrt{x}+1 = 1 \Leftrightarrow \sqrt{x} = 0 \Leftrightarrow x = 0$ (thỏa mãn)

TH2: $\sqrt{x}+1 = 3 \Leftrightarrow \sqrt{x} = 2 \Leftrightarrow x = 4$ (loại)

Vậy $x \in \{0\}$ thì $P \in \mathbb{Z}$.

d) Xét hiệu $P - 1 = \dfrac{\sqrt{x}-2}{\sqrt{x}+1} - 1 = \dfrac{\sqrt{x}-2-\sqrt{x}-1}{\sqrt{x}+1} = \dfrac{-3}{\sqrt{x}+1} < 0, \forall x \ge 0; x \ne 4$

Suy ra $P < 1$.

e) $P = \sqrt{x}-3 \Leftrightarrow \dfrac{\sqrt{x}-2}{\sqrt{x}+1} = \sqrt{x}-3 \Leftrightarrow \sqrt{x}-2 = (\sqrt{x}-3)(\sqrt{x}+1)$
\[
\Leftrightarrow \sqrt{x}-2 = x - 3\sqrt{x} + \sqrt{x} - 3 \Leftrightarrow x - 3\sqrt{x} - 1 = 0 \quad (2)
\]
Đặt $\sqrt{x} = t$ ($t \ge 0; t \ne 2$).

$(2)$ trở thành: $t^2 - 3t - 1 = 0$
\[
\Delta = (-3)^2 - 4 \cdot 1 \cdot (-1) = 13
\]
\[
\Rightarrow \hoac{&t = \dfrac{3+\sqrt{13}}{2} \text{ (TM)}\\&t = \dfrac{3-\sqrt{13}}{2} \text{ (loại)}} \Rightarrow x = \left( \dfrac{3+\sqrt{13}}{2} \right)^2 = \dfrac{22+6\sqrt{13}}{4} \text{ (thỏa mãn điều kiện)}
\]
Vậy $x = \dfrac{22+6\sqrt{13}}{4}$ là giá trị cần tìm.
}
\end{bt}

%@ id=13CD.06.020 | trang=225 | so=8 | vitri=C3.9.5,CD2.2 | muc=9 | tt=0 | ghichu=gốc tính sai nghiệm t_2 = -6 và đánh số ý 1 thay vì 2 | gl=1 | soat=sua
\begin{bt}
Cho biểu thức $A = \left( \dfrac{4\sqrt{y}}{2+\sqrt{y}} + \dfrac{8y}{4-y} \right) : \left( \dfrac{\sqrt{y}-1}{y-2\sqrt{y}} - \dfrac{2}{\sqrt{y}} \right)$, với $y > 0; y \neq 4; y \neq 9$.
\begin{cacy}
\item Rút gọn biểu thức $A$.
\item Tìm $y$ để $A = -2$.
\end{cacy}
\loigiai{
Ta có:
\[
\begin{aligned}
A &= \left( \dfrac{4\sqrt{y}\left(\sqrt{y}-2\right) - 8y}{\left(\sqrt{y}-2\right)\left(\sqrt{y}+2\right)} \right) : \left( \dfrac{\sqrt{y}-1}{\sqrt{y}\left(\sqrt{y}-2\right)} - \dfrac{2}{\sqrt{y}} \right) \\
&= \left( \dfrac{-4y-8\sqrt{y}}{\left(\sqrt{y}-2\right)\left(\sqrt{y}+2\right)} \right) : \left( \dfrac{\sqrt{y}-1-2\left(\sqrt{y}-2\right)}{\sqrt{y}\left(\sqrt{y}-2\right)} \right) \\
&= \dfrac{-4\sqrt{y}\left(\sqrt{y}+2\right)}{\left(\sqrt{y}-2\right)\left(\sqrt{y}+2\right)} : \dfrac{-\sqrt{y}+3}{\sqrt{y}\left(\sqrt{y}-2\right)} \\
&= \dfrac{-4\sqrt{y}}{\left(\sqrt{y}-2\right)} \cdot \dfrac{\sqrt{y}\left(\sqrt{y}-2\right)}{-\sqrt{y}+3} = \dfrac{4y}{\sqrt{y}-3}
\end{aligned}
\]
(với $y > 0; y \neq 4; y \neq 9$).

2) Tìm $y$ để $A = -2$:\\
Ta có $A = -2 \Leftrightarrow \dfrac{4y}{\sqrt{y}-3} = -2 \Leftrightarrow 4y = -2\left(\sqrt{y}-3\right) \Leftrightarrow 4y + 2\sqrt{y} - 6 = 0$.

Đặt $\sqrt{y} = t > 0$ ta có phương trình $4t^2 + 2t - 6 = 0$.

Ta có $a+b+c=0$ nên phương trình có hai nghiệm:
$t_1 = 1$ (thỏa mãn điều kiện)
$t_2 = -\dfrac{3}{2}$ (không thỏa mãn điều kiện)

Với $t = 1$, ta có $y = 1$ (thỏa mãn điều kiện).
Vậy $A = -2 \Leftrightarrow y = 1$.
}
\end{bt}

%@ id=13CD.06.034 | trang=232 | so=22 | vitri=C3.9.4,CD2.3 | muc=8 | tt=0 | ghichu=lời giải câu a ngắt bước tính ngoặc trước nhưng viết liền dấu bằng | gl=1 | soat=sua
\begin{bt}
Cho hai biểu thức $T = \left( \dfrac{a\sqrt{a}-1}{a-\sqrt{a}} - \dfrac{a\sqrt{a}+1}{a+\sqrt{a}} \right) : \dfrac{a+1}{a-1}$ với $a > 0; a \neq 1$.
\begin{cacy}
\item Rút gọn biểu thức $T$
\item Tìm các giá trị nguyên của $a$ để $T$ nhận giá trị nguyên
\end{cacy}
\loigiai{
\begin{cacy}
\item Ta có:
\[
T = \left( \dfrac{a\sqrt{a}-1}{a-\sqrt{a}} - \dfrac{a\sqrt{a}+1}{a+\sqrt{a}} \right) : \dfrac{a+1}{a-1} = \left( \dfrac{\left(\sqrt{a}\right)^3-1}{\left(\sqrt{a}\right)^2-\sqrt{a}} - \dfrac{\left(\sqrt{a}\right)^3+1}{\left(\sqrt{a}\right)^2+\sqrt{a}} \right) : \dfrac{a+1}{a-1} = \left(\dfrac{a+\sqrt{a}+1}{\sqrt{a}} - \dfrac{a-\sqrt{a}+1}{\sqrt{a}}\right) : \dfrac{a+1}{a-1} = 2 : \dfrac{a+1}{a-1}
\]
$\Rightarrow T = 2 : \dfrac{a+1}{a-1} = \dfrac{2(a-1)}{a+1}$.

\item Ta có $\Rightarrow T = \dfrac{2(a-1)}{a+1} = 2 - \dfrac{4}{a+1}$. Do đó $T \in \mathbb{Z}$, khi $a+1$ là ước của $4$.
Vì $a > 0, a \neq 1$ nên ta có $a+1 = 4 \Rightarrow a = 3$.
Vậy $a = 3$ là giá trị cần tìm.
\end{cacy}
}
\end{bt}

%@ id=13CD.06.046 | trang=238 | so=34 | vitri=C3.9.4,CD2.3 | muc=8 | tt=0 | gl=1
\begin{bt}
Cho biểu thức $B = \dfrac{\sqrt{a}}{\sqrt{a}-3} - \dfrac{3}{\sqrt{a}+3} - \dfrac{a-2}{a-9}$ với $a \ge 0; a \ne 9$.
\begin{cacy}
\item Rút gọn $B$.
\item Tìm các số nguyên $a$ để $B$ nhận giá trị nguyên.
\end{cacy}
\loigiai{
\begin{cacy}
\item Với $a \ge 0; a \ne 9$ ta có:
\begin{align*}
B &= \dfrac{\sqrt{a}}{\sqrt{a}-3} - \dfrac{3}{\sqrt{a}+3} - \dfrac{a-2}{\left(\sqrt{a}-3\right)\left(\sqrt{a}+3\right)} \\
&= \dfrac{\sqrt{a}\left(\sqrt{a}+3\right) - 3\left(\sqrt{a}-3\right) - (a-2)}{\left(\sqrt{a}-3\right)\left(\sqrt{a}+3\right)} \\
&= \dfrac{a+3\sqrt{a}-3\sqrt{a}+9-a+2}{\left(\sqrt{a}-3\right)\left(\sqrt{a}+3\right)} = \dfrac{11}{a-9}.
\end{align*}

\item Để $B \in \mathbb{Z} \Leftrightarrow \dfrac{11}{a-9} \in \mathbb{Z} \Leftrightarrow 11 \vdots (a-9) \Leftrightarrow (a-9) \in \text{Ư}(11)$.\\
$\text{Ư}(11) = \{1; 11; -1; -11\}$. Khi đó ta có bảng giá trị:
\begin{center}
\begin{tabular}{|c|c|c|c|c|}
\hline
$a-9$ & -11 & -1 & 1 & 11 \\ \hline
$a$ & -2 & 8 & 10 & 20 \\ \hline
& Không thỏa mãn & Thỏa mãn & Thỏa mãn & Thỏa mãn \\ \hline
\end{tabular}
\end{center}
Vậy $a \in \{8; 10; 20\}$ thì $B \in \mathbb{Z}$.
\end{cacy}
}
\end{bt}

%@ id=13CD.06.035 | trang=232 | so=23 | vitri=C3.9.4,C3.9.6,CD2.3 | muc=9 | tt=0 | gl=1
\begin{bt}
Cho biểu thức $A = \dfrac{(\sqrt{x}+1)^2+(\sqrt{x}-1)^2}{(\sqrt{x}-1)(\sqrt{x}+1)} - \dfrac{3\sqrt{x}+1}{x-1}$ với $x \ge 0, x \neq 1$.
\begin{cacy}
\item Rút gọn biểu thức $A$
\item Tìm $x$ là số chính phương để $2019A$ là số nguyên.
\end{cacy}
\loigiai{
\begin{cacy}
\item $A = \dfrac{(\sqrt{x}+1)^2 + (\sqrt{x}-1)^2 - 3\sqrt{x}-1}{x-1} = \dfrac{x+2\sqrt{x}+1+x-2\sqrt{x}+1-3\sqrt{x}-1}{x-1}$
$= \dfrac{2x-3\sqrt{x}+1}{x-1} = \dfrac{2x-2\sqrt{x}-\sqrt{x}+1}{x-1} = \dfrac{(2\sqrt{x}-1)(\sqrt{x}-1)}{(\sqrt{x}-1)(\sqrt{x}+1)} = \dfrac{2\sqrt{x}-1}{\sqrt{x}+1}$.
\item $2019A = \dfrac{2019(2\sqrt{x}+2-3)}{\sqrt{x}+1} = 4038 - \dfrac{6057}{\sqrt{x}+1}$.
$2019A$ là số nguyên khi và chỉ khi $\sqrt{x}+1$ là ước nguyên dương của 6057 gồm:
$1;3;9;673;2019;6057$.
\par +) $\sqrt{x}+1=1 \Leftrightarrow \sqrt{x}=0 \Leftrightarrow x = 0$, thỏa mãn.
\par +) $\sqrt{x}+1=3 \Leftrightarrow \sqrt{x}=2 \Leftrightarrow x = 4$, thỏa mãn.
\par +) $\sqrt{x}+1=9 \Leftrightarrow \sqrt{x}=8 \Leftrightarrow x = 64$, thỏa mãn.
\par +) $\sqrt{x}+1=673 \Leftrightarrow \sqrt{x}=672 \Leftrightarrow x = 451584$, thỏa mãn.
\par +) $\sqrt{x}+1=2019 \Leftrightarrow \sqrt{x}=2018 \Leftrightarrow x = 4072324$, thỏa mãn.
\par +) $\sqrt{x}+1=6057 \Leftrightarrow \sqrt{x}=6056 \Leftrightarrow x = 36675136$, thỏa mãn.
\end{cacy}
}
\end{bt}

%@ id=13CD.06.063 | trang=252 | so=51 | vitri=C3.9.5,CD2.3 | muc=9 | tt=0 | gl=1
\begin{bt}
Cho $A = \dfrac{x\left( \sqrt{x+4\sqrt{x-4}} + \sqrt{x-4\sqrt{x-4}} \right)}{\sqrt{x^2-8x+16}}$ với $x > 4$
\begin{cacy}
\item Rút gọn $A$. Tìm $x$ để $A$ đạt giá trị nhỏ nhất.
\item Tìm các giá trị nguyên của $x$ để $A$ có giá trị nguyên.
\end{cacy}
\loigiai{
a) Điều kiện để biểu thức $A$ xác định là $x > 4$.
\[
A = \dfrac{x\left( \sqrt{(\sqrt{x-4}+2)^2} + \sqrt{(\sqrt{x-4}-2)^2} \right)}{\sqrt{(x-4)^2}} = \dfrac{x\left( |\sqrt{x-4}+2| + |\sqrt{x-4}-2| \right)}{|x-4|}
\]
\[
= \dfrac{x(\sqrt{x-4}+2+|\sqrt{x-4}-2|)}{x-4}
\]
+ Nếu $4 < x < 8$ thì $\sqrt{x-4}-2 < 0$ nên $A = \dfrac{x(\sqrt{x-4}+2+2-\sqrt{x-4})}{x-4} = \dfrac{4x}{x-4} = 4 + \dfrac{16}{x-4}$.

Do $4 < x < 8$ nên $0 < x-4 < 4 \Rightarrow A > 8$.

+ Nếu $x \ge 8$ thì $\sqrt{x-4}-2 \ge 0$ nên
\[
A = \dfrac{x(\sqrt{x-4}+2+\sqrt{x-4}-2)}{x-4} = \dfrac{2x\sqrt{x-4}}{x-4} = \dfrac{2x}{\sqrt{x-4}} = 2\sqrt{x-4} + \dfrac{8}{\sqrt{x-4}} \ge 2\sqrt{16} = 8
\]
(Theo bất đẳng thức Cô si). Dấu bằng xảy ra khi và chỉ khi $2\sqrt{x-4} = \dfrac{8}{\sqrt{x-4}} \Leftrightarrow x-4 = 4 \Leftrightarrow x = 8$.

Vậy GTNN của $A$ bằng $8$ khi $x = 8$.

b) Xét $4 < x < 8$ thì $A = 4 + \dfrac{16}{x-4}$, ta thấy $A \in \mathbb{Z}$ khi và chỉ khi $\dfrac{16}{x-4} \in \mathbb{Z} \Leftrightarrow x-4$ là ước số nguyên dương của $16$. Hay $x-4 \in \{1; 2; 4; 8; 16\} \Leftrightarrow x \in \{5; 6; 8; 12; 20\}$ đối chiếu điều kiện suy ra $x = 5$ hoặc $x = 6$.

+ Xét $x \ge 8$ ta có: $A = \dfrac{2x}{\sqrt{x-4}}$, đặt $\sqrt{x-4} = m \Rightarrow \heva{&x = m^2+4\\&m \ge 2}$ khi đó ta có:
\[
A = \dfrac{2(m^2+4)}{m} = 2m + \dfrac{8}{m} \text{ suy ra } m \in \{2; 4; 8\} \Leftrightarrow x \in \{8; 20; 68\}.
\]
Tóm lại để $A$ nhận giá trị nguyên thì $x \in \{5; 6; 8; 20; 68\}$.
}
\end{bt}

%@ id=13CD.06.030 | trang=230 | so=18 | vitri=C3.9.4,CD2.4 | muc=8 | tt=0 | gl=1
\begin{bt}
Cho $A = \left( 1 - \dfrac{a-3\sqrt{a}}{a-9} \right) : \left( \dfrac{\sqrt{a}-2}{\sqrt{a}+3} + \dfrac{\sqrt{a}-3}{2-\sqrt{a}} - \dfrac{9-a}{a+\sqrt{a}-6} \right)$, với $a \ge 0; a \neq 4; a \neq 9$.
\begin{cacy}
\item Rút gọn $A$
\item Tìm $a$ để $A + |A| = 0$.
\end{cacy}
\loigiai{
\begin{cacy}
\item Ta có $1 - \dfrac{a-3\sqrt{a}}{a-9} = 1 - \dfrac{\sqrt{a}\left(\sqrt{a}-3\right)}{\left(\sqrt{a}-3\right)\left(\sqrt{a}+3\right)} = 1 - \dfrac{\sqrt{a}}{\sqrt{a}+3} = \dfrac{3}{\sqrt{a}+3}$.

$\dfrac{\sqrt{a}-2}{\sqrt{a}+3} + \dfrac{\sqrt{a}-3}{2-\sqrt{a}} - \dfrac{9-a}{a+\sqrt{a}-6} = \dfrac{\left(\sqrt{a}-2\right)^2 - \left(\sqrt{a}+3\right)\left(\sqrt{a}-3\right) - 9 + a}{\left(\sqrt{a}+3\right)\left(\sqrt{a}-2\right)} = \dfrac{\sqrt{a}-2}{\sqrt{a}+3}$.

$\Rightarrow A = \dfrac{3}{\sqrt{a}+3} \cdot \dfrac{\sqrt{a}+3}{\sqrt{a}-2} = \dfrac{3}{\sqrt{a}-2}$.

\item Ta có $A + |A| = 0 \Leftrightarrow |A| = -A \Leftrightarrow A \le 0 \Leftrightarrow \dfrac{3}{\sqrt{a}-2} < 0 \Leftrightarrow \sqrt{a}-2 < 0 \Leftrightarrow \sqrt{a} < 2 \Leftrightarrow 0 \le a < 4$.
Vậy $0 \le a < 4$ là giá trị cần tìm.
\end{cacy}
}
\end{bt}

%@ id=13CD.06.033 | trang=231 | so=21 | vitri=C3.9.4,CD2.4 | muc=8 | tt=0 | gl=1
\begin{bt}
Cho biểu thức: $B = \left( \dfrac{1}{3-\sqrt{x}} - \dfrac{1}{3+\sqrt{x}} \right) \cdot \dfrac{3+\sqrt{x}}{\sqrt{x}}$, (với $x > 0; x \neq 9$).
Rút gọn biểu thức và tìm tất cả các giá trị nguyên của $x$ để $B > \dfrac{1}{2}$.
\loigiai{
Ta có $B = \left( \dfrac{1}{3-\sqrt{x}} - \dfrac{1}{3+\sqrt{x}} \right) \cdot \dfrac{3+\sqrt{x}}{\sqrt{x}} = \dfrac{3+\sqrt{x} - (3-\sqrt{x})}{\left(3-\sqrt{x}\right)\left(3+\sqrt{x}\right)} \cdot \dfrac{3+\sqrt{x}}{\sqrt{x}} = \dfrac{2\sqrt{x}}{\left(3-\sqrt{x}\right)\left(3+\sqrt{x}\right)} \cdot \dfrac{3+\sqrt{x}}{\sqrt{x}} = \dfrac{2}{3-\sqrt{x}}$.

$B > \dfrac{1}{2} \Leftrightarrow \dfrac{2}{3-\sqrt{x}} > \dfrac{1}{2} \Leftrightarrow \dfrac{2}{3-\sqrt{x}} - \dfrac{1}{2} > 0 \Leftrightarrow \dfrac{4-(3-\sqrt{x})}{2\left(3-\sqrt{x}\right)} > 0 \Leftrightarrow \dfrac{1+\sqrt{x}}{2\left(3-\sqrt{x}\right)} > 0 \quad (*)$.

Vì $1+\sqrt{x} > 0$ nên $(*) \Leftrightarrow 3-\sqrt{x} > 0 \Leftrightarrow \sqrt{x} < 3 \Leftrightarrow 0 < x < 9$.
Vì $x \in \mathbb{Z} \Rightarrow x \in \{1; 2; 3; 4; 5; 6; 7; 8\}$.
}
\end{bt}

%@ id=13CD.06.045 | trang=237 | so=33 | vitri=C3.9.4,CD2.4 | muc=8 | tt=0 | gl=1 | soat=ok
\begin{bt}
Cho biểu thức $Q = \left( \dfrac{\sqrt{x}+1}{\sqrt{x}-2} - \dfrac{2\sqrt{x}}{\sqrt{x}+2} + \dfrac{5\sqrt{x}+2}{4-x} \right) : \dfrac{3\sqrt{x}-x}{x+4\sqrt{x}+4}$.
\begin{cacy}
\item Rút gọn $Q$;
\item Tìm $x$ để $Q = 2$;
\item Tìm các giá trị của $x$ để $Q$ có giá trị âm.
\end{cacy}
\loigiai{
ĐKXĐ: $x > 0; x \ne 4; x \ne 9$.
\begin{cacy}
\item 
\begin{align*}
Q &= \left( \dfrac{\sqrt{x}+1}{\sqrt{x}-2} - \dfrac{2\sqrt{x}}{\sqrt{x}+2} + \dfrac{5\sqrt{x}+2}{4-x} \right) : \dfrac{3\sqrt{x}-x}{x+4\sqrt{x}+4} \\
&= \dfrac{\left(\sqrt{x}+1\right)\left(\sqrt{x}+2\right) - 2\sqrt{x}\left(\sqrt{x}-2\right) - \left(5\sqrt{x}+2\right)}{\left(\sqrt{x}-2\right)\left(\sqrt{x}+2\right)} : \dfrac{\sqrt{x}\left(3-\sqrt{x}\right)}{\left(\sqrt{x}+2\right)^2} \\
&= \dfrac{x+3\sqrt{x}+2 - 2x + 4\sqrt{x} - 5\sqrt{x} - 2}{\left(\sqrt{x}-2\right)\left(\sqrt{x}+2\right)} \cdot \dfrac{\left(\sqrt{x}+2\right)^2}{\sqrt{x}\left(3-\sqrt{x}\right)} \\
&= \dfrac{-x+2\sqrt{x}}{\left(\sqrt{x}-2\right)\left(\sqrt{x}+2\right)} \cdot \dfrac{\left(\sqrt{x}+2\right)^2}{\sqrt{x}\left(3-\sqrt{x}\right)} \\
&= \dfrac{-\sqrt{x}\left(\sqrt{x}-2\right)}{\left(\sqrt{x}-2\right)\left(\sqrt{x}+2\right)} \cdot \dfrac{\left(\sqrt{x}+2\right)^2}{\sqrt{x}\left(3-\sqrt{x}\right)} = \dfrac{\sqrt{x}+2}{\sqrt{x}-3}.
\end{align*}

\item $Q = 2 \Leftrightarrow \dfrac{\sqrt{x}+2}{\sqrt{x}-3} = 2 \Leftrightarrow \sqrt{x}+2 = 2\sqrt{x}-6 \Leftrightarrow -\sqrt{x} = -8 \Leftrightarrow \sqrt{x} = 8 \Leftrightarrow x = 64$ (Thỏa mãn ĐKXĐ).

\item $Q < 0 \Leftrightarrow \dfrac{\sqrt{x}+2}{\sqrt{x}-3} < 0 \Leftrightarrow \sqrt{x}-3 < 0$ (vì $\sqrt{x}+2 > 0$) $\Leftrightarrow \sqrt{x} < 3 \Leftrightarrow x < 9$.\\
Kết hợp với điều kiện xác định ta có $Q < 0$ khi $0 < x < 9$ và $x \ne 4$.
\end{cacy}
}
\end{bt}

%@ id=13CD.06.042 | trang=234 | so=30 | vitri=C3.9.4,C3.9.7,CD2.5 | muc=7 | tt=0 | gl=1
\begin{bt}
Cho biểu thức $A = \left(\dfrac{\sqrt{x}-2}{x-1} - \dfrac{\sqrt{x}+2}{x+2\sqrt{x}+1}\right) \cdot \dfrac{x^2-2x+1}{2}$
\begin{cacy}
\item Rút gọn biểu thức $A$
\item Tìm $x$ để $A \ge 0$
\item Tìm giá trị lớn nhất của $A$.
\end{cacy}
\loigiai{
\begin{cacy}
\item Điều kiện $x \ge 0; x \ne 1 \Rightarrow A = -x+\sqrt{x}$
\item Ta có\\
$A \ge 0 \Leftrightarrow -x+\sqrt{x} \ge 0 \Leftrightarrow x-\sqrt{x} \le 0 \Leftrightarrow \sqrt{x}(\sqrt{x}-1) \le 0 \Leftrightarrow 0 \le \sqrt{x} \le 1 \Leftrightarrow \heva{& 0 \le x \le 1 \\& x \ge 0; x \ne 1} \Leftrightarrow 0 \le x < 1$.
\item $A = -x+\sqrt{x} = -\left(\sqrt{x}-\dfrac{1}{2}\right)^2 + \dfrac{1}{4} \le \dfrac{1}{4}, \forall x \Leftrightarrow \sqrt{x}-\dfrac{1}{2} = 0 \Leftrightarrow x = \dfrac{1}{4} \Rightarrow \text{GTLN } A = \dfrac{1}{4} \Leftrightarrow x = \dfrac{1}{4}$.
\end{cacy}
}
\end{bt}

%% doan 235-244

%@ id=13CD.06.036 | trang=233 | so=24 | vitri=C3.9.4,C3.9.5,CD2.5 | muc=8 | tt=0 | gl=1
\begin{bt}
Cho biểu thức $P = \left( \dfrac{2\sqrt{x}}{\sqrt{x}+3} + \dfrac{\sqrt{x}}{\sqrt{x}-3} - \dfrac{3x+3}{x-9} \right) : \left( \dfrac{2\sqrt{x}-2}{\sqrt{x}-3} - 1 \right)$
\begin{cacy}
\item Rút gọn $P$
\item Tìm GTNN của $P$
\end{cacy}
\loigiai{
\begin{cacy}
\item Điều kiện $x \ge 0; x \ne 9 \Rightarrow P = \dfrac{-3}{\sqrt{x}+3}$
\item Với $x \ge 0; x \ne 9 \Rightarrow \sqrt{x}+3 \ge 3 \Rightarrow \dfrac{3}{\sqrt{x}+3} \le 1 \Rightarrow P = \dfrac{-3}{\sqrt{x}+3} \ge -1$, dấu “$=$” xảy ra khi $x = 0$
Vậy $P_{\min} = -1$ khi $x = 0$.
\end{cacy}
}
\end{bt}

%@ id=13CD.06.037 | trang=233 | so=25 | vitri=C3.9.4,C3.9.6,CD2.5 | muc=8 | tt=0 | ghichu=lời giải câu b gõ nhầm \sqrt{3}+1 thay vì \sqrt{x}+1 | gl=1 | soat=sua
\begin{bt}
Cho biểu thức $A = \left(\dfrac{\sqrt{x}+3}{\sqrt{x}-2} + \dfrac{\sqrt{x}+2}{3-\sqrt{x}} + \dfrac{\sqrt{x}+2}{x-5\sqrt{x}+6}\right) : \left(1 - \dfrac{\sqrt{x}}{\sqrt{x}+1}\right)$ với $x \ge 0; x \neq 4; x \neq 9$
\begin{cacy}
\item Rút gọn biểu thức $A$
\item Với giá trị nào của $x$ thì $\dfrac{1}{A}$ đạt giá trị nhỏ nhất? Tìm giá trị nhỏ nhất đó.
\end{cacy}
\loigiai{
\begin{cacy}
\item Với $x \ge 0; x \neq 4; x \neq 9 \Rightarrow A = \dfrac{\sqrt{x}+1}{\sqrt{x}-2}$
\item Ta có $\dfrac{1}{A} = \dfrac{\sqrt{x}-2}{\sqrt{x}+1} = \dfrac{(\sqrt{x}+1)-3}{\sqrt{x}+1} = 1 - \dfrac{3}{\sqrt{x}+1}$
Với $x \ge 0; x \neq 4; x \neq 9 \Rightarrow \sqrt{x}+1 \ge 1 \Rightarrow \dfrac{-3}{\sqrt{x}+1} \ge -3 \Rightarrow \dfrac{1}{A} = 1 - \dfrac{3}{\sqrt{x}+1} \ge -2$.
Dấu “$=$” xảy ra $\Leftrightarrow x = 0$ (thỏa mãn). Khi đó $\left(\dfrac{1}{A}\right)_{\min} = -2 \Leftrightarrow x = 0$.
\end{cacy}
}
\end{bt}

%@ id=13CD.06.039 | trang=234 | so=27 | vitri=C3.9.4,C3.9.6,CD2.5 | muc=8 | tt=0 | gl=1
\begin{bt}
Cho biểu thức $P = \dfrac{1}{\sqrt{x}+1} - \dfrac{3}{x\sqrt{x}+1} + \dfrac{2}{x-\sqrt{x}+1}$
\begin{cacy}
\item Rút gọn biểu thức $P$
\item Tìm giá trị nhỏ nhất của $P$
\end{cacy}
\loigiai{
\begin{cacy}
\item Điều kiện $x \ge 0 \Rightarrow P = \dfrac{\sqrt{x}}{x-\sqrt{x}+1}$
\item Ta có $\heva{& x-\sqrt{x}+1 = \left(\sqrt{x}-\dfrac{1}{2}\right)^2 + \dfrac{3}{4} > 0, \forall x \ge 0 \\& \sqrt{x} \ge 0, \forall x \ge 0} \Rightarrow P = \dfrac{\sqrt{x}}{x-\sqrt{x}+1} \ge 0, \forall x \ge 0 \Rightarrow P_{\min} = 0 \Leftrightarrow x = 0$.
\end{cacy}
}
\end{bt}

%@ id=13CD.06.040 | trang=234 | so=28 | vitri=C3.9.4,C3.9.6,CD2.5 | muc=8 | tt=0 | gl=1 | soat=sua
\begin{bt}
Cho biểu thức $P = \dfrac{a^2+\sqrt{a}}{a-\sqrt{a}+1} - \dfrac{2a+\sqrt{a}}{\sqrt{a}} + 1$
\begin{cacy}
\item Rút gọn biểu thức $P$
\item Tìm giá trị nhỏ nhất của $P$
\end{cacy}
\loigiai{
\begin{cacy}
\item Điều kiện $a > 0 \Rightarrow P = a-\sqrt{a}$
\item Ta có $P = a-\sqrt{a} = \left(\sqrt{a}-\dfrac{1}{2}\right)^2 - \dfrac{1}{4}$. Vì $\left(\sqrt{a}-\dfrac{1}{2}\right)^2 \ge 0$ với mọi $a > 0 \Rightarrow P \ge \dfrac{-1}{4}$, dấu ``='' xảy ra khi $\left(\sqrt{a}-\dfrac{1}{2}\right)^2 = 0 \Leftrightarrow a = \dfrac{1}{4}$ (thỏa mãn).\\
Vậy $P_{\min} = \dfrac{-1}{4}$.
\end{cacy}
}
\end{bt}

%@ id=13CD.06.043 | trang=235 | so=31 | vitri=C3.9.4,CD2.5 | muc=8 | tt=0 | gl=1
\begin{bt}
Cho biểu thức $Q = \dfrac{x+2\sqrt{x}-10}{x-\sqrt{x}-6} - \dfrac{\sqrt{x}-2}{\sqrt{x}-3} - \dfrac{1}{\sqrt{x}+2}$ ($x \ge 0; x \ne 9$).
\begin{cacy}
\item Rút gọn biểu thức $Q$.
\item Tính giá trị của $Q$ khi $x = 16$.
\item Tìm giá trị của $x$ khi $Q = \dfrac{1}{3}$.
\item Tìm giá trị của $x$ sao cho $Q > \dfrac{1}{9}$.
\item Tìm giá trị lớn nhất của $Q$.
\end{cacy}
\loigiai{
\begin{cacy}
\item Với $x \ge 0; x \ne 9$ thì
\begin{align*}
Q &= \dfrac{x+2\sqrt{x}-10}{x-3\sqrt{x}+2\sqrt{x}-6} - \dfrac{\sqrt{x}-2}{\sqrt{x}-3} - \dfrac{1}{\sqrt{x}+2} \\
&= \dfrac{x+2\sqrt{x}-10}{\sqrt{x}\left(\sqrt{x}-3\right)+2\left(\sqrt{x}-3\right)} - \dfrac{\sqrt{x}-2}{\sqrt{x}-3} - \dfrac{1}{\sqrt{x}+2} \\
&= \dfrac{x+2\sqrt{x}-10}{\left(\sqrt{x}-3\right)\left(\sqrt{x}+2\right)} - \dfrac{\sqrt{x}-2}{\sqrt{x}-3} - \dfrac{1}{\sqrt{x}+2} \\
&= \dfrac{x+2\sqrt{x}-10 - \left(\sqrt{x}-2\right)\left(\sqrt{x}+2\right) - \left(\sqrt{x}-3\right)}{\left(\sqrt{x}-3\right)\left(\sqrt{x}+2\right)} \\
&= \dfrac{x+2\sqrt{x}-10 - x + 4 - \sqrt{x} + 3}{\left(\sqrt{x}-3\right)\left(\sqrt{x}+2\right)} \\
&= \dfrac{\sqrt{x}-3}{\left(\sqrt{x}-3\right)\left(\sqrt{x}+2\right)} = \dfrac{1}{\sqrt{x}+2}.
\end{align*}
Vậy với $x \ge 0; x \ne 9$ thì $Q = \dfrac{1}{\sqrt{x}+2}$.

\item Thay $x = 16$ (thỏa mãn $x \ge 0; x \ne 9$) vào $Q$ ta được:
\[ Q = \dfrac{1}{\sqrt{16}+2} = \dfrac{1}{4+2} = \dfrac{1}{6}. \]
Vậy khi $x = 16$ thì $Q = \dfrac{1}{6}$.

\item $Q = \dfrac{1}{3} \Leftrightarrow \dfrac{1}{\sqrt{x}+2} = \dfrac{1}{3} \Leftrightarrow 3 = \sqrt{x}+2 \Leftrightarrow \sqrt{x} = 1 \Leftrightarrow x = 1$ (thỏa mãn $x \ge 0; x \ne 9$).\\
Vậy với $x = 1$ thì $Q = \dfrac{1}{3}$.

\item $Q > \dfrac{1}{9} \Leftrightarrow \dfrac{1}{\sqrt{x}+2} > \dfrac{1}{9} \Leftrightarrow \dfrac{1}{\sqrt{x}+2} - \dfrac{1}{9} > 0 \Leftrightarrow \dfrac{9-\sqrt{x}-2}{9\left(\sqrt{x}+2\right)} > 0 \Leftrightarrow \dfrac{7-\sqrt{x}}{\sqrt{x}+2} > 0 \quad (1)$.\\
Vì $\sqrt{x} \ge 0$ với mọi $x \ge 0; x \ne 9$ nên $\sqrt{x}+2 > 0$ với mọi $x \ge 0; x \ne 9$.\\
$\Rightarrow (1) \Leftrightarrow 7-\sqrt{x} > 0 \Leftrightarrow \sqrt{x} < 7 \Leftrightarrow x < 49$.\\
Kết hợp với điều kiện $x \ge 0; x \ne 9$ nên $\heva{& 0 \le x < 49 \\& x \ne 9.}$\\
Vậy với $\heva{& 0 \le x < 49 \\& x \ne 9}$ thì $Q > \dfrac{1}{9}$.

\item Vì $\sqrt{x} \ge 0$ với mọi $x \ge 0; x \ne 9$ nên $\sqrt{x}+2 \ge 2$ với mọi $x \ge 0; x \ne 9$.\\
$\Rightarrow \dfrac{1}{\sqrt{x}+2} \le \dfrac{1}{2}$ với mọi $x \ge 0; x \ne 9$.\\
Vậy $Q$ đạt giá trị lớn nhất bằng $\dfrac{1}{2}$ khi $x = 0$ (thỏa mãn $x \ge 0; x \ne 9$).
\end{cacy}
}
\end{bt}

%@ id=13CD.06.044 | trang=236 | so=32 | vitri=C3.9.4,CD2.5 | muc=8 | tt=0 | gl=1
\begin{bt}
Cho biểu thức $P = \dfrac{3\sqrt{x}+2}{\sqrt{x}+1} - \dfrac{2\sqrt{x}-3}{3-\sqrt{x}} - \dfrac{3\left(3\sqrt{x}-5\right)}{x-2\sqrt{x}-3}$.
\begin{cacy}
\item Rút gọn $P$;
\item Tìm giá trị của $P$, biết $x = 4+2\sqrt{3}$;
\item Tìm giá trị nhỏ nhất của $P$.
\end{cacy}
\loigiai{
ĐKXĐ: $x \ge 0; x \ne 9$.
\begin{cacy}
\item 
\begin{align*}
P &= \dfrac{3\sqrt{x}+2}{\sqrt{x}+1} + \dfrac{2\sqrt{x}-3}{\sqrt{x}-3} - \dfrac{3\left(3\sqrt{x}-5\right)}{\left(\sqrt{x}+1\right)\left(\sqrt{x}-3\right)} \\
&= \dfrac{\left(3\sqrt{x}+2\right)\left(\sqrt{x}-3\right) + \left(2\sqrt{x}-3\right)\left(\sqrt{x}+1\right) - 3\left(3\sqrt{x}-5\right)}{\left(\sqrt{x}+1\right)\left(\sqrt{x}-3\right)} \\
&= \dfrac{3x-9\sqrt{x}+2\sqrt{x}-6 + 2x+2\sqrt{x}-3\sqrt{x}-3 - 9\sqrt{x}+15}{\left(\sqrt{x}+1\right)\left(\sqrt{x}-3\right)} \\
&= \dfrac{5x-17\sqrt{x}+6}{\left(\sqrt{x}+1\right)\left(\sqrt{x}-3\right)} = \dfrac{5x-15\sqrt{x}-2\sqrt{x}+6}{\left(\sqrt{x}+1\right)\left(\sqrt{x}-3\right)} \\
&= \dfrac{\left(5\sqrt{x}-2\right)\left(\sqrt{x}-3\right)}{\left(\sqrt{x}+1\right)\left(\sqrt{x}-3\right)} = \dfrac{5\sqrt{x}-2}{\sqrt{x}+1}.
\end{align*}

\item Ta có $x = 4+2\sqrt{3} = \left(\sqrt{3}+1\right)^2 \Rightarrow \sqrt{x} = \sqrt{3}+1$;\\
Do đó: $P = \dfrac{5\left(\sqrt{3}+1\right)-2}{\left(\sqrt{3}+1\right)+1} = \dfrac{5\sqrt{3}+3}{\sqrt{3}+2} = \dfrac{\left(5\sqrt{3}+3\right)\left(2-\sqrt{3}\right)}{\left(\sqrt{3}+2\right)\left(2-\sqrt{3}\right)} = 7\sqrt{3}-9$.

\item Ta có $P = \dfrac{5\sqrt{x}-2}{\sqrt{x}+1} = \dfrac{5\sqrt{x}+5-7}{\sqrt{x}+1} = 5 - \dfrac{7}{\sqrt{x}+1}$.\\
Vì $\dfrac{7}{\sqrt{x}+1} > 0$ nên $P$ có giá trị nhỏ nhất $\Leftrightarrow \dfrac{7}{\sqrt{x}+1}$ lớn nhất $\Leftrightarrow \sqrt{x}+1$ nhỏ nhất $\Leftrightarrow x = 0$.\\
Khi đó $\min P = 5-7 = -2$.
\end{cacy}
}
\end{bt}

%@ id=13CD.06.022 | trang=226 | so=10 | vitri=C3.9.6,CD2.5 | muc=9 | tt=0 | gl=1
\begin{bt}
Cho $P = \dfrac{\sqrt{x}}{\sqrt{x}+3} + \dfrac{2\sqrt{x}}{\sqrt{x}-3} - \dfrac{3x+9}{x-9}$, với $x \ge 0, x \neq 9$.
\begin{cacy}
\item Rút gọn $P$.
\item Tìm giá trị của $x$ để $P = \dfrac{1}{3}$.
\item Tìm giá trị lớn nhất của $P$.
\end{cacy}
\loigiai{
\begin{cacy}
\item $P = \dfrac{\sqrt{x}\left(\sqrt{x}-3\right) + 2\sqrt{x}\left(\sqrt{x}+3\right) - 3x - 9}{\left(\sqrt{x}-3\right)\left(\sqrt{x}+3\right)} = \dfrac{3\sqrt{x}-9}{\left(\sqrt{x}-3\right)\left(\sqrt{x}+3\right)} = \dfrac{3}{\sqrt{x}+3}$.
\item $P = \dfrac{1}{3} \Leftrightarrow \dfrac{3}{\sqrt{x}+3} = \dfrac{1}{3} \Rightarrow \sqrt{x}+3 = 9 \Leftrightarrow x = 36$ (thỏa mãn ĐKXĐ).
\item Với $x \ge 0, P = \dfrac{3}{\sqrt{x}+3} \le \dfrac{3}{0+3} = 1 \Rightarrow P_{\max} = 1$ khi $x = 0$ (TM).
\end{cacy}
}
\end{bt}

%@ id=13CD.06.038 | trang=233 | so=26 | vitri=C3.9.4,C3.9.6,CD2.5 | muc=9 | tt=0 | gl=1
\begin{bt}
Cho biểu thức $A = \left( \dfrac{x+2}{x\sqrt{x}-1} + \dfrac{\sqrt{x}}{x+\sqrt{x}+1} + \dfrac{1}{1-\sqrt{x}} \right) : \dfrac{\sqrt{x}-1}{2}$
\begin{cacy}
\item Rút gọn biểu thức $A$
\item Tìm $x$ để biểu thức $A$ đạt giá trị lớn nhất
\end{cacy}
\loigiai{
\begin{cacy}
\item Điều kiện $0 \le x \ne 1 \Rightarrow A = \dfrac{2}{x+\sqrt{x}+1}$
\item Ta có $A = \dfrac{2}{x+\sqrt{x}+1}$, với $0 \le x \ne 1 \Rightarrow x+\sqrt{x}+1 \ge 1 \Rightarrow A = \dfrac{2}{x+\sqrt{x}+1} \le 2$
Dấu “$=$” xảy ra $\Leftrightarrow x = 0 \Rightarrow A_{\max} = 2 \Leftrightarrow x = 0$.
\end{cacy}
}
\end{bt}

%@ id=13CD.06.041 | trang=234 | so=29 | vitri=C3.9.4,C3.9.6,CD2.5 | muc=9 | tt=0 | gl=1
\begin{bt}
Cho biểu thức $P = \left(\dfrac{\sqrt{x}}{\sqrt{x}-1} + \dfrac{\sqrt{x}}{x-1}\right) : \left(\dfrac{2}{x} - \dfrac{2-x}{x\sqrt{x}+x}\right)$
\begin{cacy}
\item Rút gọn biểu thức $P$
\item Tìm giá trị nhỏ nhất của $\sqrt{P}$
\end{cacy}
\loigiai{
\begin{cacy}
\item Điều kiện $x > 0; x \ne 1 \Rightarrow P = \dfrac{x}{\sqrt{x}-1}$
\item Với $x > 0$, nên $\sqrt{P}$ có nghĩa khi $P > 0 \Rightarrow \dfrac{x}{\sqrt{x}-1} > 0 \Rightarrow x > 1$.\\
Do $P > 0$ với mọi $x > 1 \Rightarrow \sqrt{P}_{\min} \Leftrightarrow P_{\min}$.\\
Xét $\dfrac{1}{P} = \dfrac{\sqrt{x}-1}{x} = \dfrac{1}{\sqrt{x}} - \dfrac{1}{x} = -\left(\dfrac{1}{\sqrt{x}}-\dfrac{1}{2}\right)^2 + \dfrac{1}{4} \le \dfrac{1}{4} \Rightarrow P \ge 4$, dấu ``='' xảy ra khi $x = 4$\\
$\Rightarrow P_{\min} = 4 \Rightarrow \sqrt{P}_{\min} = 2$ khi $x = 4$.
\end{cacy}
}
\end{bt}

%@ id=13CD.06.052 | trang=243 | so=40 | vitri=C3.9.4,CD2.5 | muc=9 | tt=0 | ghichu=lời giải câu d TH2 gốc ghi nhầm S=1<0 | gl=1 | soat=sua
\begin{bt}
Cho biểu thức $P = 1 : \left( \dfrac{x+2\sqrt{x}-2}{x\sqrt{x}+1} - \dfrac{\sqrt{x}-1}{x-\sqrt{x}+1} + \dfrac{1}{\sqrt{x}+1} \right)$ với $x > 0$.
\begin{cacy}
\item Rút gọn $P$
\item Tính giá trị của $P$ biết $x = 7-4\sqrt{3}$
\item Tìm $x$ để $P = 2\sqrt{x}-1$
\item Tìm $m$ để có giá trị $x$ thoả mãn $P = m$
\item Tìm giá trị nhỏ nhất của $P$
\end{cacy}
\loigiai{
\begin{cacy}
\item Với $x > 0$ thì
\begin{align*}
P &= 1 : \left[ \dfrac{x+2\sqrt{x}-2}{\left(\sqrt{x}+1\right)\left(x-\sqrt{x}+1\right)} - \dfrac{\sqrt{x}-1}{x-\sqrt{x}+1} + \dfrac{1}{\sqrt{x}+1} \right] \\
&= \dfrac{\left(\sqrt{x}+1\right)\left(x-\sqrt{x}+1\right)}{x+2\sqrt{x}-2-x+1+x-\sqrt{x}+1} = \dfrac{\left(\sqrt{x}+1\right)\left(x-\sqrt{x}+1\right)}{x+\sqrt{x}} \\
&= \dfrac{\left(\sqrt{x}+1\right)\left(x-\sqrt{x}+1\right)}{\sqrt{x}\left(\sqrt{x}+1\right)} = \dfrac{x-\sqrt{x}+1}{\sqrt{x}}.
\end{align*}
Vậy với $x > 0$ thì $P = \dfrac{x-\sqrt{x}+1}{\sqrt{x}}$.

\item Với $x = 7-4\sqrt{3} = 4 - 2 \cdot 2 \cdot \sqrt{3} + 3 = \left(2-\sqrt{3}\right)^2$ thỏa mãn điều kiện $x > 0$.\\
$\Rightarrow \sqrt{x} = \sqrt{\left(2-\sqrt{3}\right)^2} = \left|2-\sqrt{3}\right| = 2-\sqrt{3}$ (vì $2 > \sqrt{3}$)\\
$\Rightarrow P = \dfrac{x-\sqrt{x}+1}{\sqrt{x}} = \dfrac{7-4\sqrt{3}-\left(2-\sqrt{3}\right)+1}{2-\sqrt{3}} = \dfrac{6-3\sqrt{3}}{2-\sqrt{3}} = 3$.\\
Vậy với $x = 7-4\sqrt{3}$ thì $P = 3$.

\item $P = 2\sqrt{x}-1 \Leftrightarrow \dfrac{x-\sqrt{x}+1}{\sqrt{x}} = 2\sqrt{x}-1 \Leftrightarrow x-\sqrt{x}+1 = 2x-\sqrt{x} \Leftrightarrow x = 1$ (thỏa mãn $x > 0$).\\
Vậy với $x = 1$ thì $P = 2\sqrt{x}-1$.

\item $P = \dfrac{x-\sqrt{x}+1}{\sqrt{x}} = m \Leftrightarrow x-\sqrt{x}+1 = m\sqrt{x} \Leftrightarrow x-(m+1)\sqrt{x}+1 = 0 \quad (1)$.\\
Vì $1 \ne 0$ nên $(1)$ là phương trình bậc hai.\\
Đặt $t = \sqrt{x} \quad (t > 0)$.\\
$(1)$ trở thành $t^2 - (m+1)t + 1 = 0 \quad (2)$.\\
Ta có $\Delta = (m+1)^2 - 4 = m^2+2m-3 = m^2-m+3m-3 = m(m-1)+3(m-1) = (m-1)(m+3)$.\\
Phương trình $(1)$ có nghiệm $\Leftrightarrow$ Phương trình $(2)$ có nghiệm dương.\\
TH1: Phương trình $(2)$ có 2 nghiệm dương
\[ \heva{& 1 \ne 0 \\& \Delta = (m-1)(m+3) \ge 0 \\& S = m+1 > 0 \\& P = 1 > 0} \Leftrightarrow \heva{& m-1 \ge 0 \\& m+1 > 0} \Leftrightarrow m \ge 1. \]
TH2: Phương trình $(2)$ có 2 nghiệm trái dấu $\Leftrightarrow P = 1 < 0$ (vô lý) $\Rightarrow$ Loại.\\
TH3: Phương trình $(2)$ có 1 nghiệm dương và 1 nghiệm bằng 0.\\
Với $t = 0$ thay vào $(2)$ ta được $0^2 - (m+1) \cdot 0 + 1 = 0 \Leftrightarrow 1 = 0$ (vô lý) $\Rightarrow$ Loại.\\
Vậy $m \ge 1$ là giá trị cần tìm.

\item $P = \dfrac{x-\sqrt{x}+1}{\sqrt{x}} = \sqrt{x} + \dfrac{1}{\sqrt{x}} - 1$.\\
Vì $x > 0$ nên $\sqrt{x} > 0; \dfrac{1}{\sqrt{x}} > 0$.\\
Áp dụng bất đẳng thức Cô-si cho hai số dương $\sqrt{x}; \dfrac{1}{\sqrt{x}}$ ta được:
\[ \sqrt{x} + \dfrac{1}{\sqrt{x}} \ge 2\sqrt{\sqrt{x} \cdot \dfrac{1}{\sqrt{x}}} = 2. \]
Dấu "=" xảy ra $\Leftrightarrow \sqrt{x} = \dfrac{1}{\sqrt{x}} \Leftrightarrow x = 1$ (thỏa mãn $x > 0$).\\
$\Rightarrow P \ge 2-1 = 1$. Vậy giá trị nhỏ nhất của $P = 1$ khi $x = 1$.
\end{cacy}
}
\end{bt}

%% doan 245-254

%@ id=13CD.06.026 | trang=228 | so=14 | vitri=C3.9.4,CD2.6 | muc=8 | tt=0 | ghichu=lời giải dòng 3 gõ nhầm (\sqrt{x}+2) ở mẫu số | gl=1 | soat=sua
\begin{bt}
Cho biểu thức $A = \left( \dfrac{2}{\sqrt{x}-2} + \dfrac{3}{2\sqrt{x}+1} - \dfrac{5\sqrt{x}-7}{2x-3\sqrt{x}-2} \right) : \dfrac{2\sqrt{x}+3}{5x-10\sqrt{x}}$ ($x > 0, x \neq 4$).
\begin{cacy}
\item Rút gọn biểu thức $A$.
\item Tìm $x$ sao cho $A$ nhận giá trị là một số nguyên.
\end{cacy}
\loigiai{
\begin{cacy}
\item Với $x > 0, x \neq 4$ biểu thức có nghĩa ta có:
\[
\begin{aligned}
A &= \left( \dfrac{2}{\sqrt{x}-2} + \dfrac{3}{2\sqrt{x}+1} - \dfrac{5\sqrt{x}-7}{2x-3\sqrt{x}-2} \right) : \dfrac{2\sqrt{x}+3}{5x-10\sqrt{x}} \\
&= \dfrac{2\left(2\sqrt{x}+1\right) + 3\left(\sqrt{x}-2\right) - \left(5\sqrt{x}-7\right)}{\left(\sqrt{x}-2\right)\left(2\sqrt{x}+1\right)} : \dfrac{2\sqrt{x}+3}{5\sqrt{x}\left(\sqrt{x}-2\right)} \\
&= \dfrac{2\sqrt{x}+3}{\left(\sqrt{x}-2\right)\left(2\sqrt{x}+1\right)} \cdot \dfrac{5\sqrt{x}\left(\sqrt{x}-2\right)}{2\sqrt{x}+3} = \dfrac{5\sqrt{x}}{2\sqrt{x}+1}.
\end{aligned}
\]
Vậy với $x > 0, x \neq 4$ thì $A = \dfrac{5\sqrt{x}}{2\sqrt{x}+1}$.

\item Ta có $\sqrt{x} > 0, \forall x > 0, x \neq 4$ nên $A = \dfrac{5\sqrt{x}}{2\sqrt{x}+1} > 0, x > 0, x \neq 4$.
$A = \dfrac{5\sqrt{x}}{2\sqrt{x}+1} = \dfrac{5}{2} - \dfrac{5}{2\left(2\sqrt{x}+1\right)} < \dfrac{5}{2}, x > 0, x \neq 4 \Rightarrow 0 < A < \dfrac{5}{2}$, kết hợp với $A$ nhận giá trị là một số nguyên thì $A \in \{1, 2\}$.

$A = 1 \Leftrightarrow 5\sqrt{x} = 2\sqrt{x}+1 \Rightarrow \sqrt{x} = \dfrac{1}{3} \Leftrightarrow x = \dfrac{1}{9}$ thỏa mãn điều kiện.

$A = 2 \Leftrightarrow 5\sqrt{x} = 4\sqrt{x}+2 \Leftrightarrow \sqrt{x} = 2 \Leftrightarrow x = 4$ không thỏa mãn điều kiện.

Vậy với $x = \dfrac{1}{9}$ thì $A$ nhận giá trị là nguyên.
\end{cacy}
}
\end{bt}

%@ id=13CD.06.047 | trang=238 | so=35 | vitri=C3.9.4,CD2.6 | muc=9 | tt=0 | gl=1 | soat=ok
\begin{bt}
Cho biểu thức $A = \dfrac{x-2\sqrt{x}}{x\sqrt{x}-1} + \dfrac{\sqrt{x}+1}{x\sqrt{x}+x+\sqrt{x}} + \dfrac{1+2x-2\sqrt{x}}{x^2-\sqrt{x}}$ (Với $x > 0, x \ne 1$).
\begin{cacy}
\item Rút gọn biểu thức $A$.
\item Tìm $x$ để biểu thức $A$ nhận giá trị là số nguyên.
\end{cacy}
\loigiai{
\begin{cacy}
\item $A = \dfrac{\sqrt{x}+2}{x+\sqrt{x}+1}$.

\item \textbf{Cách 1:} Với $x > 0, x \ne 1 \Rightarrow x+\sqrt{x}+1 > \sqrt{x}+1 > 1$.\\
Vậy $0 < A = \dfrac{\sqrt{x}+2}{x+\sqrt{x}+1} < \dfrac{\sqrt{x}+2}{\sqrt{x}+1} = 1 + \dfrac{1}{\sqrt{x}+1} < 2$.\\
Vì $A$ nguyên nên $A = 1 \Leftrightarrow \dfrac{\sqrt{x}+2}{x+\sqrt{x}+1} = 1 \Leftrightarrow x = 1$ (Không thỏa mãn).\\
Vậy không có giá trị nguyên nào của $x$ để giá trị $A$ là một số nguyên.

\textbf{Cách 2: Dùng miền giá trị}
\[ A = \dfrac{\sqrt{x}+2}{x+\sqrt{x}+1} \Leftrightarrow Ax + (A-1)\sqrt{x} + A - 2 = 0 \quad (1) \]
Trường hợp 1: $A = 0 \Rightarrow \sqrt{x} = -2 \Rightarrow x \in \emptyset$.\\
Trường hợp 2: $A \ne 0 \Rightarrow \Delta = (A-1)^2 - 4A(A-2) = -3A^2 + 6A + 1 \ge 0 \Leftrightarrow A^2 - 2A - \dfrac{1}{3} \le 0$\\
$\Leftrightarrow A^2 - 2A + 1 \le \dfrac{4}{3} \Leftrightarrow (A-1)^2 \le \dfrac{4}{3} \Rightarrow A \in \{1; 2\}$ (do $A \in \mathbb{Z}, A > 0$).\\
Với $A = 1 \Rightarrow x = 1$ (loại).\\
Với $A = 2 \Leftrightarrow \dfrac{\sqrt{x}+2}{x+\sqrt{x}+1} = 2 \Leftrightarrow x = 0$ (loại).
\end{cacy}
}
\end{bt}

%@ id=13CD.06.048 | trang=239 | so=36 | vitri=C3.9.4,CD2.6 | muc=9 | tt=0 | ghichu=ý e lời giải gốc biến đổi sai tử số | gl=1 | soat=sua
\begin{bt}
Cho biểu thức $P = \dfrac{\sqrt{x}-3}{2-\sqrt{x}} + \dfrac{\sqrt{x}-2}{3+\sqrt{x}} - \dfrac{9-x}{x+\sqrt{x}-6}$ với $x \ge 0; x \ne 4$.
\begin{cacy}
\item Rút gọn $P$.
\item Tìm $x$ để $P = \dfrac{7}{12}$.
\item Tìm $x$ để $P > \dfrac{1}{2}$.
\item Tìm tất cả các giá trị nguyên của $x$ để $\dfrac{1}{P}$ nhận giá trị nguyên.
\item Tìm tất cả các giá trị hữu tỷ của $x$ để $P$ nhận giá trị nguyên.
\end{cacy}
\loigiai{
\begin{cacy}
\item Với $x \ge 0; x \ne 4$ thì
\begin{align*}
P &= \dfrac{\sqrt{x}-3}{2-\sqrt{x}} + \dfrac{\sqrt{x}-2}{3+\sqrt{x}} - \dfrac{9-x}{x+3\sqrt{x}-2\sqrt{x}-6} \\
&= \dfrac{\sqrt{x}-3}{2-\sqrt{x}} + \dfrac{\sqrt{x}-2}{3+\sqrt{x}} - \dfrac{9-x}{\sqrt{x}\left(\sqrt{x}+3\right)-2\left(\sqrt{x}+3\right)} \\
&= -\dfrac{\sqrt{x}-3}{\sqrt{x}-2} + \dfrac{\sqrt{x}-2}{3+\sqrt{x}} - \dfrac{9-x}{\left(\sqrt{x}+3\right)\left(\sqrt{x}-2\right)} \\
&= \dfrac{-\left(\sqrt{x}-3\right)\left(\sqrt{x}+3\right) + \left(\sqrt{x}-2\right)^2 - 9 + x}{\left(\sqrt{x}+3\right)\left(\sqrt{x}-2\right)} \\
&= \dfrac{9-x + \left(\sqrt{x}-2\right)^2 - 9 + x}{\left(\sqrt{x}+3\right)\left(\sqrt{x}-2\right)} = \dfrac{\left(\sqrt{x}-2\right)^2}{\left(\sqrt{x}+3\right)\left(\sqrt{x}-2\right)} = \dfrac{\sqrt{x}-2}{\sqrt{x}+3}.
\end{align*}
Vậy với $x \ge 0; x \ne 4$ thì $P = \dfrac{\sqrt{x}-2}{\sqrt{x}+3}$.

\item $P = \dfrac{7}{12} \Leftrightarrow \dfrac{\sqrt{x}-2}{\sqrt{x}+3} = \dfrac{7}{12} \Leftrightarrow 12\sqrt{x}-24 = 7\sqrt{x}+21 \Leftrightarrow 5\sqrt{x} = 45 \Leftrightarrow \sqrt{x} = 9 \Leftrightarrow x = 81$ (thỏa mãn $x \ge 0; x \ne 4$).\\
Vậy với $x = 81$ thì $P = \dfrac{7}{12}$.

\item $P > \dfrac{1}{2} \Leftrightarrow \dfrac{\sqrt{x}-2}{\sqrt{x}+3} > \dfrac{1}{2} \Leftrightarrow \dfrac{\sqrt{x}-2}{\sqrt{x}+3} - \dfrac{1}{2} > 0 \Leftrightarrow \dfrac{2\sqrt{x}-4-\sqrt{x}-3}{2\left(\sqrt{x}+3\right)} > 0 \Leftrightarrow \dfrac{\sqrt{x}-7}{\sqrt{x}+3} > 0 \quad (3)$.\\
Vì $\sqrt{x} \ge 0$ với mọi $x \ge 0; x \ne 4$ nên $\sqrt{x}+3 > 0$ với mọi $x \ge 0; x \ne 4$.\\
Nên $(3) \Leftrightarrow \sqrt{x}-7 > 0 \Leftrightarrow \sqrt{x} > 7 \Leftrightarrow x > 49$.\\
Kết hợp với điều kiện $x \ge 0; x \ne 4$.\\
Vậy $x > 49$ thì $P > \dfrac{1}{2}$.

\item Ta có $\dfrac{1}{P} = \dfrac{\sqrt{x}+3}{\sqrt{x}-2} = \dfrac{\sqrt{x}-2+5}{\sqrt{x}-2} = 1 + \dfrac{5}{\sqrt{x}-2}$.\\
$\dfrac{1}{P}$ nguyên $\Leftrightarrow \dfrac{5}{\sqrt{x}-2}$ nguyên $\Leftrightarrow 5 \vdots \left(\sqrt{x}-2\right) \Leftrightarrow \sqrt{x}-2$ là $\text{Ư}(5) = \{\pm 1; \pm 5\}$.\\
Lập bảng:
\begin{center}
\begin{tabular}{|c|c|c|c|c|}
\hline
$\sqrt{x}-2$ & -1 & 1 & -5 & 5 \\ \hline
$\sqrt{x}$ & 1 & 3 & -3 & 7 \\ \hline
$x$ & 1 & 9 & Loại & 49 \\ \hline
& Thỏa mãn & Thỏa mãn & & Thỏa mãn \\ \hline
\end{tabular}
\end{center}
Vậy $x \in \{1; 9; 49\}$ thì $\dfrac{1}{P}$ nguyên.

\item Ta có $P = \dfrac{\sqrt{x}-2}{\sqrt{x}+3} = \dfrac{\sqrt{x}+3-5}{\sqrt{x}+3} = 1 - \dfrac{5}{\sqrt{x}+3}$.\\
Vì $\dfrac{5}{\sqrt{x}+3} > 0$ nên $P < 1$ với mọi $x \ge 0; x \ne 4$.\\
Mà $\sqrt{x}+3 \ge 3 \Rightarrow \dfrac{5}{\sqrt{x}+3} \le \dfrac{5}{3} \Rightarrow -\dfrac{5}{\sqrt{x}+3} \ge -\dfrac{5}{3} \Rightarrow 1 - \dfrac{5}{\sqrt{x}+3} \ge 1 - \dfrac{5}{3} = -\dfrac{2}{3}$.\\
Do đó $-\dfrac{2}{3} \le P < 1$. Mà $P \in \mathbb{Z}$ nên $P = 0 \Leftrightarrow \sqrt{x} = 2 \Leftrightarrow x = 4$ (loại). Vậy không có giá trị hữu tỷ nào của $x$ để $P$ nguyên.
\end{cacy}
}
\end{bt}

%@ id=13CD.06.051 | trang=242 | so=39 | vitri=C3.9.4,CD2.6 | muc=9 | tt=0 | ghichu=lời giải câu e viết nhầm mẫu thành \sqrt{x}+3 ở các TH | gl=1 | soat=sua
\begin{bt}
Cho biểu thức $B = \dfrac{2(x+4)}{x-3\sqrt{x}-4} + \dfrac{\sqrt{x}}{\sqrt{x}+1} - \dfrac{8}{\sqrt{x}-4}$ với $x \ge 0; x \ne 16$.
\begin{cacy}
\item Rút gọn $B$.
\item Tìm giá trị của $x$ để $B = 1$.
\item Tính giá trị của $x$ sao cho $B$ không vượt quá $\dfrac{3}{2}$.
\item Tìm giá trị của $B$ khi $x$ thỏa mãn đẳng thức $\sqrt{2x-1} = x$.
\item Tìm $x$ để giá trị của $B$ là một số nguyên.
\end{cacy}
\loigiai{
\begin{cacy}
\item Với $x \ge 0; x \ne 16$ thì
\begin{align*}
B &= \dfrac{2(x+4)}{x+\sqrt{x}-4\sqrt{x}-4} + \dfrac{\sqrt{x}}{\sqrt{x}+1} - \dfrac{8}{\sqrt{x}-4} \\
&= \dfrac{2(x+4)}{\sqrt{x}\left(\sqrt{x}+1\right)-4\left(\sqrt{x}+1\right)} + \dfrac{\sqrt{x}}{\sqrt{x}+1} - \dfrac{8}{\sqrt{x}-4} \\
&= \dfrac{2(x+4)}{\left(\sqrt{x}-4\right)\left(\sqrt{x}+1\right)} + \dfrac{\sqrt{x}}{\sqrt{x}+1} - \dfrac{8}{\sqrt{x}-4} \\
&= \dfrac{2x+8+\sqrt{x}\left(\sqrt{x}-4\right)-8\left(\sqrt{x}+1\right)}{\left(\sqrt{x}-4\right)\left(\sqrt{x}+1\right)} \\
&= \dfrac{2x+8+x-4\sqrt{x}-8\sqrt{x}-8}{\left(\sqrt{x}-4\right)\left(\sqrt{x}+1\right)} = \dfrac{3x-12\sqrt{x}}{\left(\sqrt{x}-4\right)\left(\sqrt{x}+1\right)} \\
&= \dfrac{3\sqrt{x}\left(\sqrt{x}-4\right)}{\left(\sqrt{x}-4\right)\left(\sqrt{x}+1\right)} = \dfrac{3\sqrt{x}}{\sqrt{x}+1}.
\end{align*}
Vậy với $x \ge 0; x \ne 16$ thì $B = \dfrac{3\sqrt{x}}{\sqrt{x}+1}$.

\item $B = 1 \Leftrightarrow \dfrac{3\sqrt{x}}{\sqrt{x}+1} = 1 \Leftrightarrow 3\sqrt{x} = \sqrt{x}+1 \Leftrightarrow 2\sqrt{x} = 1 \Leftrightarrow \sqrt{x} = \dfrac{1}{2} \Leftrightarrow x = \dfrac{1}{4}$ (thỏa mãn $x \ge 0; x \ne 16$).\\
Vậy $x = \dfrac{1}{4}$ thì $B = 1$.

\item $B$ không vượt quá $\dfrac{3}{2} \Leftrightarrow B \le \dfrac{3}{2} \Leftrightarrow \dfrac{3\sqrt{x}}{\sqrt{x}+1} \le \dfrac{3}{2} \Leftrightarrow \dfrac{3\sqrt{x}}{\sqrt{x}+1} - \dfrac{3}{2} \le 0$\\
$\Leftrightarrow \dfrac{6\sqrt{x}-3\sqrt{x}-3}{2\left(\sqrt{x}+1\right)} \le 0 \Leftrightarrow \dfrac{3\sqrt{x}-3}{2\left(\sqrt{x}+1\right)} \le 0 \Leftrightarrow \dfrac{\sqrt{x}-1}{\sqrt{x}+1} \le 0 \quad (*)$\\
Vì $\sqrt{x} \ge 0$ với mọi $x \ge 0; x \ne 16$ nên $\sqrt{x}+1 > 0$ với mọi $x \ge 0; x \ne 16$.\\
Suy ra $(*)\Leftrightarrow \sqrt{x}-1 \le 0 \Leftrightarrow \sqrt{x} \le 1 \Leftrightarrow x \le 1$.\\
Kết hợp với điều kiện $x \ge 0; x \ne 16$.\\
Vậy $0 \le x \le 1$ thì $B$ không vượt quá $\dfrac{3}{2}$.

\item Ta có $\sqrt{2x-1} = x$ ($x \ge 0; x \ne 16$) $\Leftrightarrow 2x-1 = x^2 \Leftrightarrow x^2-2x+1 = 0 \Leftrightarrow (x-1)^2 = 0 \Leftrightarrow x = 1$ (thỏa mãn $x \ge 0; x \ne 16$).\\
$\Rightarrow B = \dfrac{3\sqrt{1}}{\sqrt{1}+1} = \dfrac{3}{2}$.\\
Vậy $\sqrt{2x-1} = x$ thì $B = \dfrac{3}{2}$.

\item $B = \dfrac{3\sqrt{x}}{\sqrt{x}+1} = \dfrac{3\sqrt{x}+3-3}{\sqrt{x}+1} = 3 - \dfrac{3}{\sqrt{x}+1} < 3$ (vì $\dfrac{3}{\sqrt{x}+1} > 0$ với mọi $x \ge 0; x \ne 16$).\\
Vì $\sqrt{x} \ge 0$ với mọi $x \ge 0; x \ne 16$ nên $\sqrt{x}+1 \ge 1$ với mọi $x \ge 0; x \ne 16$\\
$\Rightarrow \dfrac{3}{\sqrt{x}+1} \le 3 \Rightarrow -\dfrac{3}{\sqrt{x}+1} \ge -3 \Rightarrow 3 - \dfrac{3}{\sqrt{x}+1} \ge 0$.\\
Suy ra $0 \le B < 3$. Mà $B \in \mathbb{Z}$ nên $B \in \{0; 1; 2\}$.\\
TH1: $B = 0 \Leftrightarrow \dfrac{3\sqrt{x}}{\sqrt{x}+1} = 0 \Leftrightarrow x = 0$ (thỏa mãn).\\
TH2: $B = 1 \Leftrightarrow \dfrac{3\sqrt{x}}{\sqrt{x}+1} = 1 \Leftrightarrow 3\sqrt{x} = \sqrt{x}+1 \Leftrightarrow \sqrt{x} = \dfrac{1}{2} \Leftrightarrow x = \dfrac{1}{4}$ (thỏa mãn).\\
TH3: $B = 2 \Leftrightarrow \dfrac{3\sqrt{x}}{\sqrt{x}+1} = 2 \Leftrightarrow 3\sqrt{x} = 2\sqrt{x}+2 \Leftrightarrow \sqrt{x} = 2 \Leftrightarrow x = 4$ (thỏa mãn).\\
Vậy $x \in \left\{0; \dfrac{1}{4}; 4\right\}$ thì $B \in \mathbb{Z}$.
\end{cacy}
}
\end{bt}
